\documentclass{itkmitlproject}

\usepackage{afterpage}
\usepackage{graphicx,amsmath,amssymb,amsthm}
\usepackage{indentfirst}
\usepackage[hidelinks]{hyperref}
\usepackage{hyphenat}
% ระบบอ้างอิงแบบหมายเลขในเครื่องหมายก้ามปู เช่น [1], [2] และเรียงตามลำดับที่
% ปรากฏในเนื้อหา (ห้ามใช้ option sort หรือรูปแบบนาม-ปี)
\usepackage[numbers,square]{natbib}
\usepackage{longtable}          % ตารางที่ยาวข้ามหน้า
\usepackage{tikz}               % แผนภาพประกอบ
\usetikzlibrary{positioning,arrows.meta}
\usepackage{placeins}           % \FloatBarrier กันรูปไม่ให้ลอยข้ามหัวข้อ
% ให้ยืดช่องว่างได้มากขึ้นเมื่อบรรทัดมีคำภาษาอังกฤษยาวที่ตัดไม่ได้
% เพื่อไม่ให้ข้อความล้นเข้าขอบกระดาษ (คู่มือ ข้อ 4.4)
\setlength{\emergencystretch}{3em}

% ชื่อพารามิเตอร์และชื่อไฟล์ (\texttt) ใช้ TH SarabunPSK เช่นเดียวกับเนื้อหา (คู่มือ ข้อ 4.2.2)
\setmonofont{THSarabunPSK}[
	Path = fonts/,
	UprightFont = THSarabun,
	ItalicFont = THSarabun Italic,
	BoldFont = THSarabun Bold,
	BoldItalicFont = THSarabun BoldItalic,
	Extension = .ttf,
	Script=Thai
]

% ลูกศรบอกทิศของการส่งผ่าน เช่น 2017\dirarrow{}2018
%   TH SarabunPSK ไม่มีอักขระ → จึงวาดเส้นเองให้ความหนาตามน้ำหนักของตัวอักษรรอบข้าง
%   แทนการใช้ $\to$ ซึ่งดึงฟอนต์ Times มาปนกับข้อความไทย
\makeatletter
\DeclareRobustCommand{\dirarrow}{%
	\begingroup
	\def\dir@w{0.06em}%
	\ifdefstring{\f@series}{bx}{\def\dir@w{0.1em}}{}%
	\ifdefstring{\f@series}{b}{\def\dir@w{0.1em}}{}%
	\kern0.1em\tikz[baseline=-0.62ex]\draw[line width=\dir@w, line cap=round,
		-{Straight Barb[length=0.28em, width=0.52em, line cap=round]}] (0,0) -- (0.9em,0);\kern0.1em
	\endgroup}
\makeatother
\pdfstringdefDisableCommands{\def\dirarrow{→}}

% \useEnglish
\useThai

\graphicspath{ {images/} }

% ข้อความที่ยังต้องตรวจหรือเติม พิมพ์เป็นสีแดง ลบออกก่อนส่งเล่มจริง
\newcommand{\needcheck}[1]{\textcolor{red}{[ต้องตรวจอีกครั้ง: #1]}}

%1. Your thesis title (THAI)
\newcommand{\ThesisTitle}{ระบบจำแนกประเภทการโจมตีทางไซเบอร์จาก Network Logs \mbox{ด้วยปัญญาประดิษฐ์}}
%2. Your thesis title (ENG)
\newcommand{\ThesisTitleENG}{AI-Based Attack Type Classification System \mbox{from Network Logs}}
%3. Author name
\newcommand{\Author}{นายศิรเชษฐ์ โชติฐากุลอนันต์}
%4. Author name ENG
\newcommand{\AuthorENG}{Sirachet Chottakulanan}
%5. Author student ID
\newcommand{\SId}{66070191}
%6. Author 2 name
\newcommand{\AuthorTwo}{นายสุขุม ฤดีเมธากุล}
%7. Author 2 name ENG
\newcommand{\AuthorTwoENG}{Sukhum Rudeemaetakul}
%8. Author 2 student ID
\newcommand{\SIdTwo}{66070315}
%9. Advisor name
\newcommand{\Advisor}{ผศ.ดร. ประพันธ์ ปวรางกูร}
%10. Advisor name ENG
\newcommand{\AdvisorENG}{Asst.Prof.Dr. Praphan Pavarangkoon}
%11. ภาคเรียนที่
\newcommand{\Sem}{1}
%12. ปีการศึกษา พ.ศ.
\newcommand{\AcaY}{2569}
%13. ปีการศึกษา ค.ศ.
\newcommand{\AcaYAD}{2026}
%14. วันส่งรายงาน
\newcommand{\SubD}{10 พฤศจิกายน พ.ศ. 2569}
%15. Copyright year (AD)
\newcommand{\CopyrightYAD}{2026}

\begin{document}
    \frontmatter
    \lhead{}\rhead{}\chead{}\lfoot{}\cfoot{\thepage}\rfoot{}
    \makecover
    \makeinnercover
    \makeengcover
    \makecopyrightcover
    \makethesiscert
    \makeprojectcert

    % Setting margin for page numbering on frontmatter
    \newgeometry{top=1in, bottom=1in, left=1.5in, right=1in, includefoot}
    \setcounter{page}{1}
    \pagenumbering{Roman}

    \makeabstract{
        ในยุคที่เครือข่ายคอมพิวเตอร์และอินเทอร์เน็ตเข้ามามีบทบาทสำคัญต่อการดำเนินชีวิตและภาคธุรกิจ ภัยคุกคามทางไซเบอร์ในรูปแบบการโจมตีเครือข่ายก็ทวีความรุนแรงและมีความหลากหลายมากขึ้นตามไปด้วย ระบบตรวจจับการบุกรุกและงานวิจัยด้านการเรียนรู้ของเครื่องจำนวนมากมุ่งเน้นการตรวจจับแบบ Benign/Attack หรือการแจ้งเตือนเหตุการณ์ผิดปกติ ขณะที่การจำแนกประเภทการโจมตีอย่างละเอียดสามารถให้ข้อมูลเพิ่มเติมสำหรับการตอบสนองต่อเหตุการณ์ โครงงานนี้จึงพัฒนาระบบจำแนกประเภทการโจมตีทางไซเบอร์จากข้อมูล Network Flow ด้วยเทคนิคการเรียนรู้ของเครื่อง พร้อมทั้งอธิบายเหตุผลเบื้องหลังการตัดสินใจของโมเดลได้อย่างโปร่งใส โดยใช้ชุดข้อมูล CICIDS2017 จำนวน 2,497,980 แถว ครอบคลุม 9 กลุ่มประเภทการจราจร แบ่งชุดข้อมูลฝึกและทดสอบตามลำดับการจับข้อมูล (capture-order split) ในสัดส่วน 70/30 ภายในแต่ละคู่ของไฟล์การจับข้อมูลและคลาส เพื่อลดความเสี่ยงของการรั่วไหลของข้อมูล จากนั้นฝึกและเปรียบเทียบโมเดลการเรียนรู้ของเครื่องแบบดั้งเดิม 7 ชนิด ได้แก่ Logistic Regression, Random Forest, XGBoost, LightGBM, CatBoost, Multi-Layer Perceptron และ Stacking Ensemble โดย 6 ชนิดแรกปรับไฮเปอร์พารามิเตอร์ภายใต้งบประมาณการค้นหาเท่ากัน ส่วน Stacking ไม่ได้ปรับ ประเมินผลทั้งด้านการจำแนก การตรวจจับการโจมตี ผลกระทบเชิงปฏิบัติการ และประสิทธิภาพเชิงคำนวณ ร่วมกับเทคนิค SHAP (SHapley Additive exPlanations) เพื่ออธิบายผลการทำนายของโมเดล และพัฒนาเป็นเว็บแดชบอร์ดด้วย FastAPI และ Next.js สำหรับนำเสนอผลลัพธ์ ผลภายในชุดข้อมูลพบว่าโมเดล CatBoost ให้ผลลัพธ์ดีที่สุดในภาพรวม ด้วยค่า Macro F1-score เฉลี่ยครบทั้ง 9 คลาสประมาณ 0.887 และมีอัตราการแจ้งเตือนผิดพลาด (False Positive Rate) ต่ำกว่า 0.01 อย่างไรก็ตาม การทดลองเสริมที่วัดการส่งผ่านข้ามชุดข้อมูลแบบสองทิศทางระหว่าง CICIDS2017 และ CSE-CIC-IDS2018 บนโครงร่างคลาสร่วม 7 คลาส ให้ภาพที่ขัดแย้งกัน กล่าวคือ ค่า Macro F1-score สูงสุดลดลงเหลือเพียง 0.34 และ 0.44 ในสองทิศทาง เทียบกับ 0.86--0.88 ภายในชุดข้อมูลเดียวกัน และสองชุดข้อมูลแยกออกจากกันได้เกือบสมบูรณ์จากคุณลักษณะโฟลว์ (AUC ไม่ต่ำกว่า 0.999) ผลนี้ชี้ว่าประสิทธิภาพที่สูงภายในชุดข้อมูลเดียวไม่ได้รับประกันว่าแบบจำลองจะให้ประสิทธิภาพใกล้เคียงกันในสภาพแวดล้อมการจับข้อมูลที่ต่างออกไป ระบบที่พัฒนาขึ้นจึงเหมาะเป็นเครื่องมือสนับสนุนการตัดสินใจในสภาพแวดล้อมที่ใกล้เคียงกับข้อมูลฝึก และควรประเมินซ้ำด้วยข้อมูลของเครือข่ายปลายทางก่อนนำไปใช้
    }

    \makeabstracteng{
        The rapid expansion of networked systems has been accompanied by an increasing volume and diversity of cyber attacks. Many machine-learning-based intrusion detection studies focus primarily on binary benign-versus-attack detection, while fine-grained attack classification can provide additional information for incident response. This project develops a machine learning-based system to classify network-flow records into specific cyber attack categories, complemented by explainable-AI techniques that show the rationale behind each prediction. The system is built on the CICIDS2017 dataset, comprising 2,497,980 flow records across nine traffic classes. To reduce data-leakage risk, the pipeline uses a capture-order 70/30 train-test split within each (capture file, class) group and trains and compares seven classical machine learning models: Logistic Regression, Random Forest, XGBoost, LightGBM, CatBoost, Multi-Layer Perceptron, and a Stacking ensemble. The first six models are tuned under an equal search budget; Stacking is not tuned. Models are evaluated on classification performance, attack detection, operational impact, and computational cost. SHAP (SHapley Additive exPlanations) is used for model interpretability, and an interactive web dashboard is developed using FastAPI and Next.js. On this dataset, CatBoost achieves the best overall performance, with a macro F1-score of approximately 0.887 and a false-positive rate below 0.01. A supplementary bidirectional cross-dataset experiment between CICIDS2017 and CSE-CIC-IDS2018, using a shared seven-class schema, shows a sharp contrast: the best macro F1-scores fall to 0.34 and 0.44 in the two directions, compared with 0.86–0.88 within a dataset. The two datasets are also almost perfectly separable using flow features alone (AUC ≥ 0.999). These results indicate that strong within-dataset performance does not guarantee comparable performance in a different capture environment. The system is therefore suited to decision support in environments similar to its training data and should be re-evaluated on target-network data before use.
    }

    \makeack

    \newpage
    \addcontentsline{toc}{chapter}{\contentsname}
    \tableofcontents

    \newpage
    \addcontentsline{toc}{chapter}{\listtablename}
    \listoftables

    \newpage
    \addcontentsline{toc}{chapter}{\listfigurename}
    \listoffigures

    % Reset frontmatter page numbering margin, back to original margin from class file
    \restoregeometry

    \mainmatter
    \lhead{}\rhead{\thepage}\chead{}\lfoot{}\cfoot{}\rfoot{}

    \include{chapter1}
    \include{chapter2}
    \chapter{วิธีการดำเนินการวิจัย}
\label{chapter:experiment}

บทนี้อธิบายระเบียบวิธีที่ใช้ในการดำเนินโครงงาน ประกอบด้วยภาพรวมของกระบวนการทำงาน การเตรียมชุดข้อมูล การแบ่งชุดข้อมูลฝึกและชุดทดสอบ การพัฒนาแบบจำลอง เกณฑ์การประเมินผล การอธิบายผลของแบบจำลอง การพัฒนาระบบต้นแบบ และมาตรการที่ใช้ควบคุมให้การทดลองสามารถทำซ้ำได้

\section{ภาพรวมของกระบวนการดำเนินงาน}

การดำเนินงานแบ่งออกเป็น 5 ระยะที่เชื่อมต่อกันเป็นไปป์ไลน์เดียว โดยผลลัพธ์ของแต่ละระยะเป็นอินพุตของระยะถัดไป ดังแสดงในรูปที่ \ref{fig:workflow}

\begin{figure}[ht]
    \centering
    \begin{tikzpicture}[
        node distance = 7mm,
        box/.style = {
            draw, rounded corners=2pt, thick,
            text width=88mm, align=center,
            minimum height=10mm, inner sep=2mm
        },
        arr/.style = {-latex, thick}
    ]
        \node[box] (p1) {ระยะที่ 1: การศึกษาและเตรียมข้อมูล};
        \node[box, below=of p1] (p2) {ระยะที่ 2: การพัฒนาแบบจำลอง};
        \node[box, below=of p2] (p3) {ระยะที่ 3: การประเมินผลและเปรียบเทียบ};
        \node[box, below=of p3] (p4) {ระยะที่ 4: การอธิบายผลของแบบจำลอง};
        \node[box, below=of p4] (p5) {ระยะที่ 5: การแสดงผลและสรุปผล};
        \draw[arr] (p1) -- (p2);
        \draw[arr] (p2) -- (p3);
        \draw[arr] (p3) -- (p4);
        \draw[arr] (p4) -- (p5);
    \end{tikzpicture}
    \caption{ภาพรวมกระบวนการดำเนินงานของโครงงาน}
    \label{fig:workflow}
\end{figure}

\section{การศึกษาและเตรียมข้อมูล}

\subsection{ชุดข้อมูลที่ใช้}

โครงงานนี้ใช้ชุดข้อมูล CICIDS2017 รุ่น \textit{MachineLearningCVE} เป็นคอร์ปัสหลักเพียงชุดเดียว ชุดข้อมูลนี้ประกอบด้วยไฟล์บันทึกการจับข้อมูล (capture file) จำนวน 8 ไฟล์ ครอบคลุมการดำเนินการโจมตีตามตารางเวลาที่กำหนดไว้ล่วงหน้าเป็นเวลา 5 วันทำการ \cite{sharafaldin2018cicids2017} รายละเอียดของไฟล์และประเภทการโจมตีที่ปรากฏในแต่ละไฟล์แสดงในตารางที่ \ref{tab:capture-files}

\begin{table}[ht]
    \caption{ไฟล์บันทึกการจับข้อมูลในชุดข้อมูล CICIDS2017 และประเภทการโจมตีที่ปรากฏ}
    \label{tab:capture-files}
    \small
    \begin{tabularx}{\linewidth}{@{}l X l@{}}
        \hline
        \textbf{วันที่บันทึก} & \textbf{ชื่อไฟล์} & \textbf{ประเภทการโจมตี} \\
        \hline
        จันทร์ & Monday-WorkingHours & ไม่มี (BENIGN) \\
        อังคาร & Tuesday-WorkingHours & Brute Force \\
        พุธ & Wednesday-workingHours & DoS, Heartbleed \\
        พฤหัสบดี (เช้า) & Thursday-WorkingHours-Morning-WebAttacks & Web Attack \\
        พฤหัสบดี (บ่าย) & Thursday-WorkingHours-Afternoon-Infilteration & Infiltration \\
        ศุกร์ (เช้า) & Friday-WorkingHours-Morning & Bot \\
        ศุกร์ (บ่าย) & Friday-WorkingHours-Afternoon-PortScan & PortScan \\
        ศุกร์ (บ่าย) & Friday-WorkingHours-Afternoon-DDos & DDoS \\
        \hline
    \end{tabularx}
\end{table}

การใช้ไฟล์ทั้ง 8 ไฟล์เป็นเงื่อนไขบังคับ เนื่องจากการโจมตีแต่ละตระกูลปรากฏอยู่ในไฟล์เพียงไฟล์เดียว หากไฟล์ใดขาดหายไปจะทำให้คลาสนั้นหายไปจากชุดข้อมูลทั้งคลาส กระบวนการโหลดข้อมูลจึงได้รับการออกแบบให้ตรวจสอบความครบถ้วนของไฟล์และหยุดการทำงานพร้อมแจ้งข้อผิดพลาด แทนที่จะสร้างแคชข้อมูลที่ไม่สมบูรณ์

\subsection{การสำรวจข้อมูลเบื้องต้น}

ก่อนการเตรียมข้อมูล คณะผู้จัดทำได้ดำเนินการสำรวจข้อมูลเบื้องต้น (Exploratory Data Analysis) เพื่อทำความเข้าใจโครงสร้างของชุดข้อมูล ประกอบด้วยการตรวจสอบการกระจายตัวของคลาส การตรวจสอบค่าว่างและค่าอนันต์ในแต่ละคอลัมน์ การวิเคราะห์สหสัมพันธ์ระหว่างคุณลักษณะ และการพิจารณาการกระจายตัวของคุณลักษณะเชิงตัวเลข ผลการสำรวจยืนยันความไม่สมดุลอย่างรุนแรงของคลาส ซึ่งเป็นข้อมูลนำเข้าสำคัญต่อการออกแบบกลยุทธ์การจัดการความไม่สมดุลในระยะถัดไป

\subsection{การทำความสะอาดข้อมูล}

ขั้นตอนการทำความสะอาดข้อมูลประกอบด้วยการดำเนินการต่อไปนี้

\begin{enumerate}
    \item ตัดอักขระช่องว่างส่วนเกินออกจากชื่อคอลัมน์ทั้งหมด เนื่องจากไฟล์ต้นฉบับมีชื่อคอลัมน์ที่ขึ้นต้นด้วยช่องว่าง
    \item ลบคอลัมน์ Fwd Header Length.1 ซึ่งเป็นคอลัมน์ที่ซ้ำซ้อนกับคอลัมน์ Fwd Header Length และเป็นข้อบกพร่องที่ทราบแล้วของชุดข้อมูลรุ่นที่เผยแพร่
    \item ลบระเบียนที่มีค่าว่าง (missing value) และค่าอนันต์ (infinity) ในคุณลักษณะที่ใช้เป็นอินพุตของแบบจำลอง
    \item ลบระเบียนที่มีป้ายกำกับที่ไม่สามารถจับคู่เข้ากับคลาสมาตรฐานทั้ง 9 คลาสได้
    \item บันทึกตำแหน่งของแถวในไฟล์ต้นฉบับไว้ในคอลัมน์เมทาดาทา \_row\_index ก่อนการลบระเบียนใด ๆ เพื่อใช้เป็นกุญแจลำดับการจับข้อมูลในขั้นตอนการแบ่งข้อมูล
\end{enumerate}

หลังการทำความสะอาด ชุดข้อมูลที่ใช้งานประกอบด้วยระเบียนโฟลว์จำนวน 2,497,980 ระเบียน โดยมีคุณลักษณะเชิงตัวเลขที่ใช้เป็นอินพุตของแบบจำลองจำนวน 77 คุณลักษณะ ส่วนคอลัมน์เมทาดาทา ได้แก่ \text{dataset\_id}, \text{source\_file}, \text{capture\_window} และ \text{\_row\_index} ถูกกันไว้ไม่ให้เข้าสู่แบบจำลอง เพื่อป้องกันการรั่วไหลของสารสนเทศที่บ่งชี้คลาสโดยตรง

\subsection{การปรับป้ายกำกับให้เป็นมาตรฐาน}

ป้ายกำกับดิบในชุดข้อมูล CICIDS2017 มีทั้งสิ้น 15 ค่า ซึ่งบางค่ามีความละเอียดในระดับเครื่องมือที่ใช้โจมตีมากกว่าระดับตระกูลของการโจมตี และบางค่ามีอักขระควบคุมที่ไม่ใช่รหัสแอสกีปนอยู่ในสายอักขระ กระบวนการปรับป้ายกำกับให้เป็นมาตรฐานจึงดำเนินการตามลำดับดังนี้ ทำการนอร์มัลไลซ์ยูนิโคดแบบ NFKD แทนที่อักขระควบคุมด้วยช่องว่างหนึ่งตัว รวบช่องว่างที่ติดกันให้เหลือตัวเดียว แปลงเป็นตัวพิมพ์เล็ก และตัดช่องว่างหัวท้ายออก ผลลัพธ์ที่ได้ถูกนำไปจับคู่เข้ากับคลาสมาตรฐาน 9 คลาสตามตารางที่ \ref{tab:label-mapping}

\begin{table}[ht]
    \caption{การจับคู่ป้ายกำกับดิบในชุดข้อมูล CICIDS2017 เข้ากับคลาสมาตรฐาน 9 คลาส}
    \label{tab:label-mapping}
    \small
    \begin{tabularx}{\linewidth}{@{}X l@{}}
        \hline
        \textbf{ป้ายกำกับดิบ} & \textbf{คลาสมาตรฐาน} \\
        \hline
        BENIGN & BENIGN \\
        DoS Hulk, DoS GoldenEye, DoS slowloris, DoS Slowhttptest & DoS \\
        DDoS & DDoS \\
        PortScan & PortScan \\
        FTP-Patator, SSH-Patator & Brute Force \\
        Web Attack -- Brute Force, Web Attack -- XSS, Web Attack -- Sql Injection & Web Attack \\
        Bot & Bot \\
        Infiltration & Infiltration \\
        Heartbleed & Heartbleed \\
        \hline
    \end{tabularx}
\end{table}

เพื่อรองรับป้ายกำกับรูปแบบใหม่ที่อาจเพิ่มเข้ามาในอนาคต กระบวนการจับคู่ยังมีตารางสำรองที่พิจารณาจากคำนำหน้าของป้ายกำกับ เช่น ป้ายกำกับที่เริ่มต้นด้วย \texttt{dos} จะถูกจับคู่เข้ากับคลาส DoS โดยอัตโนมัติ จากนั้นป้ายกำกับที่เป็นข้อความจะถูกแปลงเป็นจำนวนเต็มด้วยการเข้ารหัสป้ายกำกับ (label encoding) เพื่อให้แบบจำลองนำไปประมวลผลได้

\subsection{การจัดการคุณลักษณะ}

การปรับสเกลของคุณลักษณะดำเนินการด้วยตัวปรับมาตรฐาน (standard scaler) ซึ่งแปลงคุณลักษณะแต่ละตัวให้มีค่าเฉลี่ยเป็นศูนย์และส่วนเบี่ยงเบนมาตรฐานเป็นหนึ่ง ตามสมการ

\begin{equation}
    z_{j} = \frac{x_{j} - \mu_{j}}{\sigma_{j}}
\end{equation}

โดยที่ $\mu_{j}$ และ $\sigma_{j}$ คือค่าเฉลี่ยและส่วนเบี่ยงเบนมาตรฐานของคุณลักษณะที่ $j$ ซึ่งประมาณค่าจากชุดฝึกเท่านั้น

ประเด็นสำคัญเชิงระเบียบวิธีคือ ตัวปรับสเกลถูกประกอบเข้าเป็นขั้นตอนแรกของไปป์ไลน์ร่วมกับตัวจำแนกประเภท มิได้ปรับสเกลข้อมูลไว้ล่วงหน้าแยกจากกระบวนการฝึก ด้วยวิธีนี้ ในการตรวจสอบไขว้และการค้นหาไฮเปอร์พารามิเตอร์ทุกครั้ง ค่าสถิติของตัวปรับสเกลจะถูกประมาณใหม่จากข้อมูลในโฟลด์ฝึกเท่านั้น จึงไม่เกิดการรั่วไหลของข้อมูลจากโฟลด์ตรวจสอบเข้าสู่กระบวนการฝึก \cite{pedregosa2011scikit}

สำหรับการคัดเลือกคุณลักษณะ โครงงานนี้เลือกใช้คุณลักษณะทั้ง 77 ตัวเป็นอินพุตของแบบจำลอง แทนการคัดกรองล่วงหน้า ด้วยเหตุผลสองประการ ประการแรก แบบจำลองที่อาศัยต้นไม้ตัดสินใจมีกลไกคัดเลือกคุณลักษณะอยู่ในตัวผ่านการเลือกจุดแบ่ง ประการที่สอง การคัดกรองคุณลักษณะโดยพิจารณาความสัมพันธ์กับป้ายกำกับจากทั้งชุดข้อมูลก่อนการแบ่งข้อมูลเป็นรูปแบบหนึ่งของการรั่วไหลของข้อมูลที่พบได้บ่อย \cite{arp2022dosdonts} การวิเคราะห์ความสำคัญของคุณลักษณะจึงกระทำภายหลังการฝึกด้วยเทคนิค SHAP ตามที่กล่าวถึงในหัวข้อ \ref{sec:shap}

\section{การแบ่งชุดข้อมูลฝึกและชุดทดสอบ}
\label{sec:split}

\subsection{หลักการแบ่งข้อมูลตามลำดับการจับข้อมูล}

โครงงานนี้ใช้การแบ่งข้อมูลตามลำดับการจับข้อมูล (capture-order split) ในสัดส่วน 70/30 ภายในทุกกลุ่มของคู่ (ไฟล์การจับข้อมูล, คลาส) กล่าวคือ ภายในแต่ละกลุ่ม ระเบียน 70\% แรกตามลำดับแถวในไฟล์ถูกจัดเป็นชุดฝึก และระเบียน 30\% ท้ายถูกจัดเป็นชุดทดสอบ เงื่อนไขนี้ทำให้ภายในกลุ่มเดียวกัน ไม่มีระเบียนในชุดทดสอบใดอยู่ก่อนระเบียนในชุดฝึกตามลำดับการจับข้อมูล และลดโอกาสที่ระเบียนซึ่งอยู่ติดกันและมีลักษณะใกล้เคียงกันจะกระจายไปอยู่ทั้งสองฝั่ง อย่างไรก็ตาม การแบ่งนี้ไม่ได้ตัดที่จุดเวลาเดียวกันทั้งคอร์ปัส จึงไม่ได้จำลองสถานการณ์ที่ฝึกด้วยข้อมูลในอดีตทั้งหมดแล้วทำนายข้อมูลในอนาคตอย่างสมบูรณ์

เนื่องจากชุดข้อมูลรุ่น \textit{MachineLearningCVE} ไม่มีคอลัมน์เวลา กุญแจลำดับที่ใช้จึงเป็นตำแหน่งของแถวในไฟล์ต้นฉบับ มิใช่เวลาประทับ (timestamp) ลำดับนี้ได้รับการตรวจสอบความสอดคล้องกับตารางเวลาการโจมตีที่ผู้จัดทำชุดข้อมูลเผยแพร่ไว้ทุกครั้งก่อนการฝึก ข้อควรระวังที่ต้องบันทึกไว้คือ เครื่องมือ CICFlowMeter สร้างระเบียนเมื่อโฟลว์สิ้นสุดลง ลำดับแถวจึงน่าจะสะท้อนลำดับที่โฟลว์สิ้นสุด ไม่ใช่ลำดับที่โฟลว์เริ่มต้น ซึ่งสอดคล้องกับระบบตรวจจับการบุกรุกที่ทำงานกับคุณลักษณะระดับโฟลว์ เพราะคุณลักษณะเหล่านั้นคำนวณได้เมื่อโฟลว์สิ้นสุดแล้วเท่านั้น

\subsection{ข้อจำกัดที่ทำให้ไม่สามารถกันข้อมูลทั้งไฟล์ได้}

การกันข้อมูลทั้งไฟล์ไว้เป็นชุดทดสอบ (source holdout) ซึ่งเป็นแนวทางที่เข้มงวดกว่า ไม่สามารถทำได้บนคอร์ปัสนี้ เพราะการโจมตีแต่ละคลาสปรากฏอยู่ในไฟล์เดียว การกันไฟล์หนึ่งออกไปจะทำให้คลาสในไฟล์นั้นไม่มีตัวอย่างในชุดฝึกเลย ด้วยเหตุผลเดียวกัน การตรวจสอบไขว้จึงใช้วิธี Stratified $k$-Fold แทน Group $k$-Fold ที่จัดกลุ่มตามไฟล์ต้นทาง ข้อจำกัดนี้ถูกบันทึกไว้อย่างชัดเจนเพื่อประกอบการตีความผลลัพธ์

\subsection{จำนวนตัวอย่างในแต่ละคลาส}

ผลของการแบ่งข้อมูลตามหลักการข้างต้นทำให้ได้ชุดฝึกจำนวน 1,748,579 ระเบียน และชุดทดสอบจำนวน 749,401 ระเบียน โดยมีจำนวนตัวอย่างในแต่ละคลาสตามตารางที่ \ref{tab:class-split}

\begin{table}[ht]
    \caption{จำนวนตัวอย่างในชุดฝึกและชุดทดสอบจำแนกตามคลาส}
    \label{tab:class-split}
    \small
    \begin{tabularx}{\linewidth}{@{}X r r r@{}}
        \hline
        \textbf{คลาส} & \textbf{ชุดฝึก} & \textbf{ชุดทดสอบ} & \textbf{รวม} \\
        \hline
        BENIGN & 1,450,574 & 621,680 & 2,072,254 \\
        DoS & 135,611 & 58,119 & 193,730 \\
        DDoS & 89,609 & 38,405 & 128,014 \\
        PortScan & 63,485 & 27,209 & 90,694 \\
        Brute Force & 6,405 & 2,745 & 9,150 \\
        Web Attack & 1,500 & 643 & 2,143 \\
        Bot & 1,363 & 585 & 1,948 \\
        Infiltration & 25 & 11 & 36 \\
        Heartbleed & 7 & 4 & 11 \\
        \hline
        \textbf{รวม} & \textbf{1,748,579} & \textbf{749,401} & \textbf{2,497,980} \\
        \hline
    \end{tabularx}
\end{table}

จากตารางที่ \ref{tab:class-split} เห็นได้ว่าคลาส BENIGN คิดเป็นสัดส่วนประมาณ 83.0\% ของชุดทดสอบ ซึ่งหมายความว่าแบบจำลองที่ทำนายว่าเป็น BENIGN ทุกกรณีจะได้ค่าความแม่นยำโดยรวมสูงถึง 0.8296 โดยไม่ตรวจจับการโจมตีได้เลยแม้แต่รายการเดียว ค่าดังกล่าวจึงถูกใช้เป็นแบบจำลองฐานเชิงเปรียบเทียบ (majority-class baseline) ในการรายงานผลทุกครั้ง นอกจากนี้ คลาส Heartbleed และ Infiltration มีจำนวนตัวอย่างในชุดทดสอบเพียงหลักหน่วยและหลักสิบตามลำดับ ค่าตัวชี้วัดรายคลาสของทั้งสองคลาสจึงต้องตีความอย่างระมัดระวังในฐานะค่าประมาณขอบเขตบน มิใช่ค่าประมาณที่เสถียร

\section{การพัฒนาแบบจำลอง}

\subsection{การกำหนดปัญหา}

ปัญหาถูกกำหนดให้เป็นการจำแนกประเภทแบบหลายคลาสที่มีจำนวนคลาส $K = 9$ โดยแบบจำลองรับเวกเตอร์คุณลักษณะ $\mathbf{x} \in \mathbb{R}^{77}$ และให้ผลลัพธ์เป็นเวกเตอร์ความน่าจะเป็นเหนือคลาสทั้งเก้า คลาสที่ถูกทำนายคือคลาสที่มีค่าความน่าจะเป็นสูงสุด ตามสมการ

\begin{equation}
    \hat{y} = \arg\max_{k \in \{1, \ldots, K\}} P(y = k \mid \mathbf{x})
\end{equation}

\subsection{แบบจำลองที่นำมาเปรียบเทียบ}

แบบจำลองที่นำมาเปรียบเทียบมี 7 รูปแบบ ครอบคลุมทั้งแบบจำลองเชิงเส้น แบบจำลองรวมกลุ่มชนิด Bagging และ Boosting โครงข่ายประสาทเทียม และการรวมกลุ่มเชิงซ้อน ตามที่สรุปไว้ในตารางที่ \ref{tab:model-summary} สำหรับ Stacking Ensemble โครงสร้างที่ใช้ประกอบด้วยแบบจำลองฐานสามตัว ได้แก่ LightGBM, XGBoost และ Random Forest โดยใช้ค่าความน่าจะเป็นของคลาสจากการทำนายนอกโฟลด์ (out-of-fold prediction) ด้วยการตรวจสอบไขว้ 5 โฟลด์ เป็นคุณลักษณะระดับเมทา และใช้ Logistic Regression เป็นแบบจำลองระดับที่สอง \cite{wolpert1992stacked, ting1999issues}

\subsection{การจัดการความไม่สมดุลของคลาส}

โครงงานนี้ใช้การถ่วงน้ำหนักคลาสแบบผกผันตามความถี่เป็นกลยุทธ์หลักและใช้แนวคิดเดียวกันในทุกแบบจำลอง เพื่อให้การจัดการความไม่สมดุลเป็นตัวแปรควบคุม โดยแบบจำลองส่วนใหญ่ (Logistic Regression, LightGBM, XGBoost และ Multi-Layer Perceptron) คำนวณน้ำหนักของคลาส $k$ จากสัดส่วนคลาสทั้งชุดฝึกครั้งเดียวตามสูตรมาตรฐานของ scikit-learn

\begin{equation}
    w_{k} = \frac{n}{K \cdot n_{k}}
\end{equation}

โดยที่ $n$ คือจำนวนตัวอย่างทั้งหมดในชุดฝึก $K$ คือจำนวนคลาส และ $n_{k}$ คือจำนวนตัวอย่างของคลาส $k$ ในชุดฝึก น้ำหนักนี้ถูกนำไปคูณกับพจน์ความสูญเสียของตัวอย่างในคลาสนั้น ๆ ทำให้ความผิดพลาดในคลาสส่วนน้อยมีต้นทุนสูงขึ้นตามสัดส่วน

โมเดลอีกสองรูปแบบใช้สูตรที่แปรผันตรงกับ $w_k$ ข้างต้นแต่คำนวณด้วยกลไกเฉพาะของแต่ละไลบรารี ได้แก่ Random Forest ใช้โหมด balanced\_subsample ซึ่งคำนวณน้ำหนักคลาสใหม่จากสัดส่วนคลาสภายในแต่ละตัวอย่างบูตสแตรป (bootstrap sample) ของต้นไม้แต่ละต้น แทนที่จะคำนวณครั้งเดียวจากทั้งชุดฝึก และ CatBoost ใช้โหมด auto\_class\_weights="Balanced" ซึ่งคำนวณจาก $w_k = \max_j(n_j) / n_k$ อันเป็นสูตรที่แปรผันตรงกับ $w_k$ ข้างต้นแต่ต่างค่าคงที่การนอร์มัลไลซ์ ทั้งสองกรณียังคงเป็นการถ่วงน้ำหนักแบบผกผันตามความถี่ในเชิงสัดส่วนระหว่างคลาสเหมือนกัน จึงไม่กระทบต่อการเปรียบเทียบระหว่างแบบจำลอง แต่ค่าคงที่ในการคำนวณมีความแตกต่างไปจากสมการข้างต้นตามข้อจำกัดของ API ในแต่ละไลบรารี

ประเด็นที่ต้องเน้นย้ำคือ การถ่วงน้ำหนักและการปรับสัดส่วนตัวอย่างใด ๆ กระทำเฉพาะภายในชุดฝึกและภายในโฟลด์ของการตรวจสอบไขว้เท่านั้น การกระจายตัวของคลาสในชุดทดสอบไม่ถูกแก้ไข เพื่อให้ผลการประเมินสะท้อนอัตราพื้นฐานที่เป็นจริง ซึ่งเป็นเงื่อนไขจำเป็นต่อการตีความอัตราการแจ้งเตือนผิดพลาดอย่างถูกต้อง \cite{axelsson2000baserate}

\subsection{การปรับไฮเปอร์พารามิเตอร์ภายใต้งบประมาณเดียวกัน}

เพื่อให้การเปรียบเทียบระหว่างแบบจำลองสะท้อนความสามารถของตัวแบบจำลองมากกว่าความพยายามในการปรับแต่ง โครงงานนี้กำหนดเงื่อนไขการค้นหาไฮเปอร์พารามิเตอร์ให้เหมือนกันใน 6 แบบจำลองที่ผ่านการค้นหา ได้แก่ Logistic Regression, Random Forest, XGBoost, LightGBM, CatBoost และ Multi-Layer Perceptron ดังนี้

\begin{itemize}
    \item วิธีการค้นหา: การค้นหาแบบสุ่ม (Randomized Search) ร่วมกับการตรวจสอบไขว้แบบแบ่งชั้น 5 โฟลด์
    \item ฟังก์ชันวัตถุประสงค์: ค่า Macro F1-score
    \item งบประมาณการค้นหา: 20 รอบการทดลองต่อแบบจำลอง บนตัวอย่างย่อยของชุดฝึกจำนวน 200,000 ระเบียน
    \item ขนาดปริภูมิการค้นหา: 144 ชุดค่าสำหรับทั้ง 6 แบบจำลอง ซึ่งทำให้สัดส่วนของปริภูมิที่ถูกสำรวจเท่ากันที่ประมาณ 13.9\%
\end{itemize}

การกำหนดขนาดปริภูมิการค้นหาให้เท่ากันมีความสำคัญ เพราะการใช้จำนวนรอบการทดลองเท่ากันบนปริภูมิที่มีขนาดต่างกันย่อมทำให้แบบจำลองที่มีปริภูมิเล็กกว่าได้รับการสำรวจในสัดส่วนที่สูงกว่า อย่างไรก็ตาม เงื่อนไขที่เท่ากันนี้ครอบคลุมเพียง 6 แบบจำลองข้างต้น Stacking Ensemble ไม่ได้ผ่านการค้นหาไฮเปอร์พารามิเตอร์ ทั้งในระดับแบบจำลองฐานและแบบจำลองระดับที่สอง เนื่องจากการค้นหาแต่ละรอบต้องฝึกแบบจำลองฐานทั้งชุดใหม่ทั้งหมด ทำให้ต้นทุนสูงเกินกว่าทรัพยากรที่มี แบบจำลองฐานของ Stacking จึงใช้ค่าตั้งต้นที่กำหนดไว้ในโค้ด มิใช่ค่าที่ได้จากการค้นหาในตารางที่ \ref{tab:best-params} ผลของ Stacking Ensemble จึงไม่ได้อยู่ภายใต้เงื่อนไขการปรับแต่งเดียวกับอีก 6 แบบจำลอง และควรตีความแยกจากการเปรียบเทียบภายใต้งบประมาณเท่ากัน ค่าไฮเปอร์พารามิเตอร์ที่ดีที่สุดของแต่ละแบบจำลองแสดงในตารางที่ \ref{tab:best-params}

\begin{table}[ht]
    \caption{ค่าไฮเปอร์พารามิเตอร์ที่ดีที่สุดของแต่ละแบบจำลองจากการค้นหาแบบสุ่ม}
    \label{tab:best-params}
    \small
    \begin{tabularx}{\linewidth}{@{}l >{\raggedright\arraybackslash}X@{}}
        \hline
        \textbf{แบบจำลอง} & \textbf{ค่าไฮเปอร์พารามิเตอร์ที่ดีที่สุด} \\
        \hline
        Logistic Regression & {C}~=~300.0, \texttt{tol}~=~$1\times10^{-4}$, \texttt{fit\_intercept}~=~True \\
        Random Forest & \texttt{n\_estimators}~=~300, \texttt{max\_depth}~=~30, \texttt{min\_samples\_split}~=~2, \texttt{min\_samples\_leaf}~=~1 \\
        XGBoost & \texttt{n\_estimators}~=~800, \texttt{max\_depth}~=~10, \texttt{learning\_rate}~=~0.05, \texttt{subsample}~=~0.8 \\
        LightGBM & \texttt{n\_estimators}~=~1000, \texttt{num\_leaves}~=~127, \texttt{learning\_rate}~=~0.1, \texttt{min\_child\_samples}~=~1 \\
        CatBoost & \texttt{iterations}~=~500, \texttt{depth}~=~6, \texttt{learning\_rate}~=~0.15, \texttt{l2\_leaf\_reg}~=~1 \\
        Multi-Layer Perceptron & \texttt{hidden\_layer\_sizes}~=~(256,~128,~64,~32), \texttt{learning\_rate\_init}~=~$5\times10^{-4}$, \texttt{batch\_size}~=~256, \texttt{alpha}~=~$1\times10^{-4}$ \\
        Stacking Ensemble & ไม่ได้ปรับไฮเปอร์พารามิเตอร์ (ใช้ค่าตั้งต้นในโค้ด) \\
        \hline
    \end{tabularx}
\end{table}

\section{การประเมินผลและเปรียบเทียบ}

\subsection{ตัวชี้วัดที่ใช้}

กำหนดให้ $TP_{k}$, $FP_{k}$, $TN_{k}$ และ $FN_{k}$ เป็นจำนวนตัวอย่างที่ทำนายถูกในคลาสบวก ทำนายผิดเป็นคลาสบวก ทำนายถูกในคลาสลบ และทำนายผิดเป็นคลาสลบ ของคลาส $k$ ตามลำดับ ค่าความเที่ยง ค่าความไว และค่า F1-score ของคลาส $k$ คำนวณได้จาก

\begin{equation}
    \mathrm{Precision}_{k} = \frac{TP_{k}}{TP_{k} + FP_{k}}, \qquad
    \mathrm{Recall}_{k} = \frac{TP_{k}}{TP_{k} + FN_{k}}
\end{equation}

\begin{equation}
    \mathrm{F1}_{k} = 2 \cdot \frac{\mathrm{Precision}_{k} \cdot \mathrm{Recall}_{k}}{\mathrm{Precision}_{k} + \mathrm{Recall}_{k}}
\end{equation}

ค่า Macro F1-score ซึ่งเป็นตัวชี้วัดหลักของโครงงานนี้ คือค่าเฉลี่ยเลขคณิตของ F1-score รายคลาสโดยให้น้ำหนักเท่ากันทุกคลาส ตามสมการ

\begin{equation}
    \mathrm{F1_{macro}} = \frac{1}{K} \sum_{k=1}^{K} \mathrm{F1}_{k}
\end{equation}

ในการทดลองหลักกำหนด $K=9$ และใช้ Macro F1 ของทั้ง 9 คลาสเป็นเกณฑ์จัดอันดับหลัก นอกจากนี้รายงาน Macro F1 เฉพาะคลาสที่มีตัวอย่างทดสอบไม่น้อยกว่า 30 ระเบียนเป็นตัวชี้วัดประกอบ ซึ่งในชุดทดสอบนี้เหลือ 7 คลาสหลังไม่นับ Heartbleed และ Infiltration ในค่าเฉลี่ย ส่วนการทดลองเสริมใช้ $K=7$ ตามโครงร่างคลาสร่วมของสองชุดข้อมูล ซึ่งรวม Infiltration แต่ไม่มี PortScan และ Heartbleed ดังนั้นตัวชี้วัด 7 คลาสทั้งสองกรณีจึงใช้คนละชุดคลาส

อัตราการแจ้งเตือนผิดพลาดและอัตราการพลาดตรวจจับคำนวณได้จาก

\begin{equation}
    \mathrm{FPR}_{k} = \frac{FP_{k}}{FP_{k} + TN_{k}}, \qquad
    \mathrm{FNR}_{k} = \frac{FN_{k}}{FN_{k} + TP_{k}}
\end{equation}

และค่าสัมประสิทธิ์สหสัมพันธ์ของแมทธิวส์สำหรับกรณีทวิภาค คำนวณได้จาก

\begin{equation}
    \mathrm{MCC} = \frac{TP \cdot TN - FP \cdot FN}{\sqrt{(TP+FP)(TP+FN)(TN+FP)(TN+FN)}}
\end{equation}

ซึ่งมีค่าอยู่ในช่วง $-1$ ถึง $1$ และทนทานต่อความไม่สมดุลของคลาสมากกว่าค่าความแม่นยำโดยรวม \cite{matthews1975comparison, chicco2020mcc}

\subsection{การประเมินผลสี่มิติ}

การประเมินแบบจำลองดำเนินการใน 4 มิติ เพื่อให้ครอบคลุมประเด็นที่เกี่ยวข้องกับการนำไปใช้งาน ดังนี้

\begin{enumerate}
    \item \textbf{ประสิทธิภาพการจำแนกประเภท (Classification Performance)} วัดด้วยค่า Accuracy, Balanced Accuracy, Precision, Recall และ F1-score ทั้งแบบ macro และ weighted รวมถึงค่า MCC โดยใช้ Macro F1-score เป็นเกณฑ์ตัดสินหลัก
    \item \textbf{ความสามารถในการตรวจจับการโจมตี (Attack Detection Ability)} ยุบผลการทำนายเป็นกรณีทวิภาคระหว่างการจราจรปกติและการโจมตี แล้ววัดอัตราการตรวจจับ ค่า FPR ค่า FNR และค่าความไวต่อการตรวจจับรายคลาส
    \item \textbf{ผลกระทบเชิงปฏิบัติการ (Operational Impact)} แปลงอัตราความผิดพลาดให้เป็นจำนวนสัมบูรณ์ ได้แก่ จำนวนการแจ้งเตือนผิดพลาดต่อทราฟฟิกปกติหนึ่งหมื่นระเบียน และจำนวนโฟลว์การโจมตีที่พลาดการตรวจจับ เพื่อสะท้อนภาระที่อาจเกิดขึ้นกับผู้ปฏิบัติงาน
    \item \textbf{ประสิทธิภาพเชิงคำนวณ (Computational Efficiency)} วัดเวลาที่ใช้ในการฝึก เวลาแฝงของการทำนายที่ควอนไทล์ที่ 50 และ 95 อัตราการประมวลผล และขนาดของไฟล์แบบจำลอง โดยการวัดเวลาแฝงดำเนินการบนหน่วยประมวลผลกลางสำหรับทุกแบบจำลองเป็นชุดข้อมูลขนาด 1,000 โฟลว์ ทำซ้ำ 30 ครั้งจากตำแหน่งเริ่มต้นแบบสุ่ม และตัดผลของรอบอุ่นเครื่องรอบแรกออก
\end{enumerate}

การวัดค่าเวลาแฝงที่ควอนไทล์ที่ 95 มีความสำคัญเทียบเท่าค่ามัชฌิม เนื่องจากระบบตรวจจับที่ทำงานเร็วเป็นส่วนใหญ่แต่ช้าเป็นครั้งคราวยังคงมีโอกาสสูญเสียโฟลว์ที่ประมวลผลไม่ทัน ซึ่งในทางปฏิบัติเทียบเท่ากับการพลาดตรวจจับ

\subsection{เกณฑ์การจัดอันดับแบบจำลอง}

เนื่องจากไม่มีตัวชี้วัดใดที่สามารถตัดสินความเหมาะสมของแบบจำลองได้ครบทุกมิติ โครงงานนี้จึงประกาศเกณฑ์การจัดอันดับไว้สามชุดควบคู่กัน แทนการเลือกแบบจำลองที่ดีที่สุดเพียงตัวเดียว ดังนี้

\begin{enumerate}
    \item \textbf{อันดับภาพรวม (Overall Best)} ตัดสินด้วยค่า Macro F1-score สูงสุด และใช้ค่า Weighted F1-score เป็นเกณฑ์ตัดสินกรณีเสมอ
    \item \textbf{อันดับด้านความมั่นคงปลอดภัย (Security-focused Best)} ตัดสินด้วยค่าความไวต่อการตรวจจับเฉลี่ยแบบ macro สูงสุด เฉพาะในกลุ่มแบบจำลองที่มีค่า FPR แบบทวิภาคไม่เกิน 0.01 และไม่มีคลาสใดที่มีค่าความไวเป็นศูนย์
    \item \textbf{อันดับด้านการนำไปใช้งาน (Deployment Best)} ตัดสินด้วยค่าเวลาแฝงที่ควอนไทล์ที่ 95 ต่ำสุด เฉพาะในกลุ่มแบบจำลองที่มีค่า Macro F1-score ต่างจากแบบจำลองที่ดีที่สุดไม่เกิน 0.02 และใช้ขนาดของไฟล์แบบจำลองเป็นเกณฑ์ตัดสินกรณีเสมอ
\end{enumerate}

เกณฑ์ทั้งสามชุดถูกระบุไว้ก่อนการวิเคราะห์ผล (criteria specified before analysis) ในไฟล์นโยบายการจัดอันดับ \texttt{configs/ranking\_policy.json} และนำไปใช้กับทุกแบบจำลองด้วยกฎเดียวกัน เพื่อลดโอกาสการเลือกเกณฑ์ให้เข้ากับผลลัพธ์ภายหลัง

\subsection{การตรวจสอบความน่าเชื่อถือของผลลัพธ์}

เพื่อตรวจว่าค่าตัวชี้วัดที่สูงอาจเกิดจากการรั่วไหลของข้อมูลหรือความบังเอิญของการแบ่งข้อมูลเพียงครั้งเดียว โครงงานนี้กำหนดการตรวจสอบเพิ่มเติมสามรายการ ประกอบด้วย การเปรียบเทียบกับแบบจำลองฐานที่ทำนายคลาสส่วนใหญ่เสมอ การรายงานค่าเฉลี่ยและส่วนเบี่ยงเบนมาตรฐานของ Macro F1-score จากการตรวจสอบไขว้ 5 โฟลด์บนชุดฝึก และการรายงานค่า Macro F1-score เฉพาะกลุ่มคลาสที่มีจำนวนตัวอย่างในชุดทดสอบไม่น้อยกว่า 30 ระเบียน ควบคู่กับค่า Macro F1-score ของทุกคลาส เพื่อให้ผู้อ่านแยกแยะได้ว่าค่าตัวชี้วัดส่วนใดได้รับผลกระทบจากคลาสที่มีตัวอย่างน้อยเกินกว่าจะประมาณค่าได้อย่างเสถียร

\section{การอธิบายผลของแบบจำลอง}
\label{sec:shap}

การวิเคราะห์เพื่ออธิบายผลของแบบจำลองดำเนินการด้วยเทคนิค SHAP ตามหลักการที่กล่าวถึงในหัวข้อ \ref{sec:xai-review} โดยใช้อัลกอริทึม TreeSHAP สำหรับแบบจำลองที่อาศัยต้นไม้ตัดสินใจ ได้แก่ Random Forest, XGBoost, LightGBM และ CatBoost ซึ่งให้ผลลัพธ์ที่แม่นยำในเวลาแบบพหุนาม \cite{lundberg2020treeshap} และใช้ Kernel SHAP ซึ่งเป็นวิธีประมาณค่าแบบไม่ขึ้นกับชนิดของแบบจำลองสำหรับ Multi-Layer Perceptron โดยคำนวณบนตัวอย่างย่อยเนื่องจากมีต้นทุนการคำนวณสูง

ผลการวิเคราะห์ถูกนำเสนอในสองระดับ ระดับแรกคือระดับภาพรวม (global explanation) ซึ่งจัดอันดับคุณลักษณะตามค่าเฉลี่ยของค่าสัมบูรณ์ของค่า SHAP เพื่อระบุว่าคุณลักษณะใดมีอิทธิพลต่อการตัดสินใจของแบบจำลองมากที่สุดโดยรวม และระดับที่สองคือระดับรายตัวอย่าง (local explanation) ซึ่งแสดงว่าคุณลักษณะแต่ละตัวผลักผลการทำนายของโฟลว์หนึ่งระเบียนไปในทิศทางใดและด้วยขนาดเท่าใด นอกจากการอธิบายผลแล้ว การวิเคราะห์นี้ยังทำหน้าที่ตรวจสอบด้วยว่าแบบจำลองเรียนรู้จากคุณลักษณะที่สอดคล้องกับลักษณะทางเทคนิคของการโจมตีจริง มากกว่าจากลักษณะเฉพาะของกระบวนการสร้างชุดข้อมูล

\section{การแสดงผลและระบบต้นแบบ}

ผลลัพธ์ทั้งหมดถูกนำเสนอผ่านช่องทางต่อไปนี้

\begin{enumerate}
    \item ตารางเปรียบเทียบตัวชี้วัดของแบบจำลองทั้ง 7 รูปแบบในทั้งสี่มิติของการประเมิน
    \item เมทริกซ์ความสับสน (Confusion Matrix) แบบจำนวนนับของแบบจำลองทั้ง 7 รูปแบบจากการทดลองหลักบน CICIDS2017 โดยแสดง CatBoost ในหัวข้อ \ref{sec:main-classification} และอีก 6 แบบจำลองในภาคผนวก \ref{appendix:main-confusion-matrices}
    \item แผนภาพสรุปผลการวิเคราะห์ด้วย SHAP ทั้งในระดับภาพรวมและระดับรายตัวอย่าง
    \item ระบบต้นแบบในรูปเว็บแดชบอร์ด พัฒนาด้วยเฟรมเวิร์ก FastAPI สำหรับส่วนบริการเบื้องหลัง และ Next.js สำหรับส่วนติดต่อผู้ใช้ ซึ่งรองรับการเรียกดูผลการเปรียบเทียบแบบจำลอง การตรวจสอบการกระจายตัวของข้อมูล การเรียกดูผลการอธิบายแบบจำลอง และการอัปโหลดไฟล์ข้อมูลโฟลว์ใหม่เพื่อทำนายผลเป็นชุด
\end{enumerate}

ส่วนของการทำนายผลข้อมูลใหม่มีการตรวจสอบโครงสร้างของไฟล์ข้อมูลนำเข้า (schema validation) โดยเปรียบเทียบกับรายชื่อคุณลักษณะที่บันทึกไว้พร้อมกับไฟล์แบบจำลอง หากไฟล์ข้อมูลนำเข้าขาดคุณลักษณะที่จำเป็น ระบบจะปฏิเสธการทำนายและแจ้งรายชื่อคุณลักษณะที่ขาดไป เพื่อป้องกันการให้ผลลัพธ์ที่ไม่ถูกต้องจากข้อมูลนำเข้าที่ไม่สอดคล้องกับที่แบบจำลองได้รับการฝึกไว้

\section{การทดลองเสริม: การส่งผ่านข้ามชุดข้อมูลแบบสองทิศทาง}
\label{sec:crossdataset}

\subsection{คำถามของการทดลองและเหตุผลของการออกแบบ}

การประเมินผลทั้งหมดในหัวข้อก่อนหน้ากระทำอยู่ภายในชุดข้อมูล CICIDS2017 ซึ่งบันทึกจากสภาพแวดล้อมเครือข่ายจำลองเพียงชุดเดียว แม้การแบ่งข้อมูลตามลำดับการจับข้อมูลจะช่วยลดความเสี่ยงของการรั่วไหลของข้อมูลภายในคอร์ปัส แต่ยังไม่สามารถตอบได้ว่าแบบจำลองเรียนรู้พฤติกรรมของการโจมตีที่ใช้ได้ทั่วไป หรือเรียนรู้ลักษณะเฉพาะของสภาพแวดล้อมการจับข้อมูลนั้น โครงงานนี้จึงจัดการทดลองเสริมขึ้นเพื่อตอบคำถามว่า \emph{แบบจำลองที่ฝึกจากคอร์ปัสหนึ่งยังคงใช้งานได้ดีเพียงใดเมื่อนำไปทำนายข้อมูลจากอีกคอร์ปัสหนึ่งที่บันทึกจากสภาพแวดล้อมต่างกัน} โดยใช้ชุดข้อมูล CSE-CIC-IDS2018 เป็นคอร์ปัสคู่เปรียบเทียบ

ค่า Macro F1-score ที่วัดได้จากคอร์ปัสอื่นเพียงค่าเดียวไม่สามารถตีความได้โดยลำพัง เพราะไม่ทราบว่าค่าที่ต่ำนั้นเกิดจากการส่งผ่านที่ล้มเหลว หรือเกิดจากงานจำแนกประเภทบนคอร์ปัสนั้นยากกว่าโดยตัวมันเอง การตีความจึงต้องอาศัยค่าอ้างอิงสองค่าบนชุดทดสอบเดียวกัน ได้แก่ \textbf{ค่าพื้น} (floor) ซึ่งเป็นค่าของแบบจำลองฐานที่ทำนายคลาสส่วนใหญ่เสมอ และ \textbf{ค่าเพดาน} (ceiling) ซึ่งเป็นค่าที่แบบจำลองชนิดเดียวกันทำได้เมื่อฝึกและทดสอบภายในคอร์ปัสปลายทางเอง ด้วยเหตุนี้การทดลองจึงต้องฝึกแบบจำลองจากทั้งสองคอร์ปัส มิใช่ใช้ CSE-CIC-IDS2018 เป็นเพียงชุดทดสอบภายนอก ผังการทดลองจึงมีลักษณะเป็นเมทริกซ์ 2×2 ตามตารางที่ \ref{tab:crossdataset-design}

\begin{table}[ht]
    \caption{ผังการทดลองแบบสองทิศทางและบทบาทของผลการวัดในแต่ละช่อง}
    \label{tab:crossdataset-design}
    \small
    \begin{tabularx}{\linewidth}{@{}l l X@{}}
        \hline
        \textbf{คอร์ปัสที่ใช้ฝึก} & \textbf{คอร์ปัสที่ใช้ทดสอบ} & \textbf{บทบาทของผลการวัด} \\
        \hline
        CICIDS2017 & CICIDS2017 & ค่าเพดานของทิศทาง 2018\dirarrow{}2017 \\
        CICIDS2017 & CSE-CIC-IDS2018 & การส่งผ่านทิศทาง 2017\dirarrow{}2018 \\
        CSE-CIC-IDS2018 & CSE-CIC-IDS2018 & ค่าเพดานของทิศทาง 2017\dirarrow{}2018 \\
        CSE-CIC-IDS2018 & CICIDS2017 & การส่งผ่านทิศทาง 2018\dirarrow{}2017 \\
        \hline
    \end{tabularx}
\end{table}

แบบจำลองในการทดลองเสริมนี้ถูกฝึกขึ้นใหม่แยกจากแบบจำลองหลักของโครงงาน โดยใช้โครงร่างคลาสร่วม 7 คลาส ตัวอย่างย่อย และค่าไฮเปอร์พารามิเตอร์ที่ไม่ผ่านการปรับแต่ง ค่าตัวชี้วัดในหัวข้อนี้จึงไม่สามารถนำไปเทียบตรงกับผลการทดลองหลักแบบ 9 คลาสได้ และไม่มีผลต่อการคัดเลือกแบบจำลองแชมเปียนที่รายงานในบทที่ \ref{chapter:result}

\subsection{ชุดข้อมูลและการปรับให้เทียบกันได้}

ชุดข้อมูล CICIDS2017 หลังการทำความสะอาดมี 2,497,980 ระเบียน (หัวข้อ \ref{sec:split}) ส่วนชุดข้อมูล CSE-CIC-IDS2018 \cite{cic2018web} ประกอบด้วยไฟล์บันทึกการจับข้อมูลรายวัน 10 ไฟล์ รวม 16,232,943 ระเบียน และเหลือ 11,445,955 ระเบียนหลังการทำความสะอาดด้วยกฎชุดเดียวกัน เนื่องจากทั้งสองคอร์ปัสมีชื่อคอลัมน์คุณลักษณะและชุดป้ายกำกับไม่ตรงกัน จึงต้องปรับให้อยู่ในรูปแบบร่วมกัน (canonicalization) ตามลำดับดังนี้

\begin{enumerate}
    \item \textbf{การจับคู่คุณลักษณะ} คุณลักษณะร่วม 77 ตัวถูกจับคู่ตามความหมายของค่าที่ CICFlowMeter สร้าง มิใช่การเทียบชื่อคอลัมน์ เนื่องจาก 50 จาก 77 คอลัมน์มีชื่อต่างกันระหว่างสองคอร์ปัส ความถูกต้องของการจับคู่ได้รับการตรวจสอบด้วยความสัมพันธ์เชิงพีชคณิตระหว่างคอลัมน์ สเกลของค่า และชนิดข้อมูล ไม่พบคู่ใดที่บ่งชี้ว่าจับคู่ผิด แม้บางคู่จะมีสเกลค่าต่างกันมาก (เช่น Fwd IAT Mean ต่างกันประมาณ 4.5 อันดับของเลขยกกำลังสิบ)
    \item \textbf{การรวมคลาส} ป้ายกำกับของทั้งสองคอร์ปัสถูกรวบเข้าเป็นคลาสร่วม 7 คลาส ได้แก่ Benign, DDoS, DoS, Brute Force, Bot, Infiltration และ Web Attack คลาส PortScan และ Heartbleed ของ CICIDS2017 ถูกตัดออกเพราะไม่มีคลาสคู่กันใน CSE-CIC-IDS2018 ทำให้ CICIDS2017 เหลือ 2,407,275 ระเบียน
    \item \textbf{ลำดับขั้นตอน} ทั้งสองคอร์ปัสผ่านขั้นตอนเดียวกันตามลำดับ คือ เปลี่ยนชื่อคอลัมน์เป็นรูปแบบร่วม รวมคลาสและตัดแถวที่ไม่มีคลาสคู่กัน ทำความสะอาด แล้วจึงสุ่มตัวอย่าง การทำความสะอาดก่อนการสุ่มทำให้ตัวอย่างทุกแถวที่สุ่มได้เป็นแถวที่ใช้งานได้จริงทั้งสองฝั่ง
    \item \textbf{การสุ่มตัวอย่าง} สุ่มตัวอย่างคอร์ปัสละ 300,000 ระเบียนด้วยค่าเริ่มต้นการสุ่ม 42 โดยเก็บทุกแถวของคลาสที่มีน้อยกว่า 10,000 ระเบียน แล้วแบ่งจำนวนที่เหลือให้คลาสใหญ่ตามสัดส่วนเดิมของคอร์ปัส วิธีนี้รักษาทั้งคลาสหายากและสัดส่วนความไม่สมดุลเดิมไว้ (Benign ประมาณ 83--88\% ทั้งสองฝั่ง) ตัวอย่างชุดนี้สุ่มครั้งเดียวและใช้ร่วมกันในทุกรอบของการทดลอง
\end{enumerate}

ตัวอย่างฝั่ง CICIDS2017 ประกอบด้วย Benign 248,188 ระเบียน DoS 23,203 ระเบียน DDoS 15,332 ระเบียน Brute Force 9,150 ระเบียน Web Attack 2,143 ระเบียน Bot 1,948 ระเบียน และ Infiltration 36 ระเบียน ขณะที่ฝั่ง CSE-CIC-IDS2018 ประกอบด้วย Benign 263,858 ระเบียน DDoS 20,272 ระเบียน DoS 5,132 ระเบียน Bot 3,778 ระเบียน Infiltration 3,642 ระเบียน Brute Force 2,459 ระเบียน และ Web Attack 859 ระเบียน คลาส Infiltration ซึ่งมีตัวอย่างเพียง 36 ระเบียนใน CICIDS2017 กลับมีตัวอย่างหลักพันใน CSE-CIC-IDS2018 ความแตกต่างของความถี่รายคลาสนี้เป็นปัจจัยที่ต้องนำมาประกอบการตีความผล

\subsection{การแบ่งชุดฝึกและชุดทดสอบ}
\label{sec:cd-split}

โพรโทคอลสุดท้ายของการทดลองเสริม (เรียกว่า protocol\_v3) แบ่งตัวอย่าง 300,000 ระเบียนของแต่ละคอร์ปัสเป็นชุดฝึก 70\% และชุดทดสอบ 30\% โดยดำเนินการสองรูปแบบคู่ขนานกัน

\begin{enumerate}
    \item \textbf{การแบ่งแบบสุ่ม (random split)} แบ่งแบบแบ่งชั้นตามคลาส (stratified) โดยสลับลำดับแถวภายในแต่ละคลาสด้วยค่าเริ่มต้นการสุ่มของรอบนั้น
    \item \textbf{การแบ่งตามลำดับการจับข้อมูล (capture-order split)} ภายในแต่ละกลุ่ม (กลุ่มการจับข้อมูล, คลาส) เรียงแถวตามกุญแจลำดับ แล้วจัด 70\% แรกเป็นชุดฝึกและ 30\% ท้ายเป็นชุดทดสอบ สำหรับ CICIDS2017 กลุ่มการจับข้อมูลคือไฟล์ต้นฉบับ (source\_file) และกุญแจลำดับคือตำแหน่งของแถวในไฟล์ CSV ต้นฉบับ (\_row\_index) เนื่องจากชุดข้อมูลรุ่น \textit{MachineLearningCVE} ไม่มีคอลัมน์เวลา ส่วน CSE-CIC-IDS2018 กลุ่มการจับข้อมูลคือวันที่บันทึก และกุญแจลำดับคือคอลัมน์ Timestamp
\end{enumerate}

ต้องระบุให้ชัดว่า การแบ่งตามลำดับการจับข้อมูลของ CICIDS2017 \emph{ไม่ใช่การแบ่งตามเวลาประทับ (timestamp split)} แต่เป็นการแบ่งตามลำดับแถวในไฟล์ ซึ่งได้รับการตรวจสอบแล้วว่าสอดคล้องกับตารางเวลาการโจมตีที่ผู้จัดทำชุดข้อมูลเผยแพร่ และการแบ่งเกิดขึ้นภายในแต่ละกลุ่ม (ไฟล์, คลาส) มิใช่การตัดที่เวลาเดียวกันทั้งคอร์ปัส เพราะการโจมตีแต่ละคลาสปรากฏอยู่ในไฟล์เดียว การตัดที่เวลาเดียวทั้งคอร์ปัสจะทำให้บางคลาสไปอยู่ฝั่งเดียวทั้งหมด ผลที่ตามมาคือ ภายในคลาสที่ประกอบด้วยเครื่องมือโจมตีหลายตัวที่ทำงานต่อเนื่องกัน ตำแหน่งที่ตัดจะกำหนดว่าเครื่องมือใดอยู่ในชุดฝึก ซึ่งส่งผลต่อผลลัพธ์อย่างมาก (หัวข้อ \ref{sec:cd-ssh})

การแบ่งตามลำดับการจับข้อมูลเป็นฟังก์ชันของข้อมูลล้วน ๆ จึงได้ชุดฝึกและชุดทดสอบเดิมทุกรอบ ส่วนการแบ่งแบบสุ่มเปลี่ยนตามค่าเริ่มต้นการสุ่ม ในโพรโทคอลก่อนหน้า (protocol\_v2) สัดส่วนการแบ่งคือ 60/10/30 โดยส่วน 10\% ถูกกันไว้สำหรับการปรับเทียบแต่ไม่ได้ถูกใช้งาน โครงงานจึงตัดส่วนดังกล่าวออกและใช้ 70/30 เป็นโพรโทคอลสุดท้าย ผลของ protocol\_v2 ถูกอ้างถึงในบทที่ \ref{chapter:result} เฉพาะเพื่อวิเคราะห์ความไวต่อสัดส่วนการแบ่งเท่านั้น

\subsection{แบบจำลอง รอบการทดลอง และตัวชี้วัด}

ในแต่ละรอบของการทดลอง (กำหนดโดยรูปแบบการแบ่งข้อมูลและค่าเริ่มต้นของการสุ่ม) แบบจำลองแต่ละรูปแบบจะถูกฝึกจากชุดฝึกของ CICIDS2017 และจากชุดฝึกของ CSE-CIC-IDS2018 แยกกัน แล้วนำแต่ละชุดไปทดสอบกับชุดทดสอบของ \emph{ทั้งสอง} คอร์ปัส ตัวปรับสเกลคุณลักษณะถูกประมาณค่าจากคอร์ปัสต้นทางเพียงฝั่งเดียวแล้วนำไปใช้กับชุดทดสอบของอีกฝั่งโดยไม่ปรับค่าใหม่ ซึ่งตรงกับกรณีที่นำแบบจำลองที่ฝึกเสร็จแล้วไปใช้กับข้อมูลของอีกสภาพแวดล้อมหนึ่งโดยไม่มีข้อมูลของสภาพแวดล้อมนั้นมาปรับเทียบ เงื่อนไขของการทดลองมีดังนี้

\begin{itemize}
    \item แบบจำลอง 7 รูปแบบเดียวกับการทดลองหลัก ได้แก่ Logistic Regression, Random Forest, XGBoost, LightGBM, CatBoost, Multi-Layer Perceptron และ Stacking Ensemble ฝึกขึ้นใหม่จากศูนย์ด้วยค่าไฮเปอร์พารามิเตอร์ตั้งต้นชุดเดียวกันทั้งสองทิศ พร้อมการถ่วงน้ำหนักคลาสแบบผกผันตามความถี่ โดยจงใจไม่นำค่าที่ปรับแต่งไว้แล้วบน CICIDS2017 (ตารางที่ \ref{tab:best-params}) มาใช้ เพราะจะทำให้ฝั่ง CICIDS2017 ได้เปรียบ
    \item ค่าเริ่มต้นการสุ่ม 5 ค่า (42--46) กับการแบ่งข้อมูล 2 รูปแบบ ได้ผลทั้งสิ้น 280 ผลการวัด (2 รูปแบบ × 5 ค่าเริ่มต้น × 2 คอร์ปัสฝึก × 2 คอร์ปัสทดสอบ × 7 แบบจำลอง)
    \item ภายใต้การแบ่งตามลำดับการจับข้อมูล ชุดฝึกและชุดทดสอบเหมือนกันทุกรอบ ส่วนเบี่ยงเบนมาตรฐานระหว่าง 5 รอบจึงวัดเฉพาะความสุ่มภายในตัวแบบจำลอง มิใช่ความผันผวนของการแบ่งข้อมูล
\end{itemize}

ตัวชี้วัดหลักคือ Macro F1-score ของทั้ง 7 คลาส รายงานเป็นค่าเฉลี่ย ± ส่วนเบี่ยงเบนมาตรฐานจาก 5 รอบ ค่าพื้นของแบบจำลองที่ทำนาย Benign ทั้งหมดเท่ากับ 0.1294 บนชุดทดสอบของ CICIDS2017 และ 0.1337 บนชุดทดสอบของ CSE-CIC-IDS2018 เพื่อให้การส่งผ่านทั้งสองทิศเทียบกันได้ทั้งที่ค่าเพดานและค่าพื้นของแต่ละคอร์ปัสไม่เท่ากัน จึงนิยามสัดส่วนประสิทธิภาพที่กู้คืนได้ (Recovery) ดังสมการ

\begin{equation}
    \mathrm{Recovery} = \frac{\mathrm{F1_{transfer}} - \mathrm{F1_{floor}}}{\mathrm{F1_{ceiling}} - \mathrm{F1_{floor}}} \times 100\%
\end{equation}

โดยที่ F1\textsubscript{transfer} คือค่าที่แบบจำลองซึ่งฝึกจากคอร์ปัสต้นทางทำได้บนชุดทดสอบของคอร์ปัสปลายทาง F1\textsubscript{floor} คือค่าพื้นบนชุดทดสอบปลายทาง และ F1\textsubscript{ceiling} คือค่าที่แบบจำลองชนิดเดียวกันทำได้เมื่อฝึกและทดสอบภายในคอร์ปัสปลายทาง ภายใต้รูปแบบการแบ่งข้อมูลเดียวกัน ค่า 100\% หมายถึงการส่งผ่านที่ไม่สูญเสียประสิทธิภาพเลย และ 0\% หมายถึงดีไม่ต่างจากการเดาคลาสส่วนใหญ่ เนื่องจากค่าเพดานแตกต่างกันในแต่ละแบบจำลอง ค่า Recovery จึงต้องอ่านคู่กับค่า Macro F1-score ดิบเสมอ และไม่ใช้จัดอันดับแบบจำลองโดยลำพัง ความไม่สมมาตรของการส่งผ่าน (asymmetry) นิยามเป็นผลต่างของค่า Recovery เฉลี่ย 7 แบบจำลองระหว่างทิศ 2018\dirarrow{}2017 กับทิศ 2017\dirarrow{}2018 มีหน่วยเป็นจุดเปอร์เซ็นต์ นอกจากนี้ยังรายงานค่า F1 รายคลาสเพื่อระบุว่าคะแนนที่เปลี่ยนไปมาจากคลาสใด

\subsection{การตรวจสอบร่องรอยเฉพาะของชุดข้อมูล}
\label{sec:cd-av-method}

เพื่อตรวจว่าสองคอร์ปัสแยกออกจากกันได้ด้วยค่าของคุณลักษณะมากเพียงใด (dataset fingerprint) โครงงานใช้การตรวจสอบเชิงปฏิปักษ์ (adversarial validation) โดยฝึกตัวจำแนก LightGBM ให้ทายว่าระเบียนหนึ่งมาจากคอร์ปัสใด แยกทีละคลาสร่วม ใช้ตัวอย่างไม่เกิน 8,000 ระเบียนต่อฝั่ง ทำซ้ำ 5 ค่าเริ่มต้นการสุ่ม และวัดด้วยค่า ROC AUC ผลดังกล่าวมีการตรวจควบคุมประกอบดังนี้

\begin{itemize}
    \item \textbf{ลำดับความจุของแบบจำลอง (model ladder)} ทำซ้ำด้วย decision stump (กฎเดียวบนคอลัมน์เดียว) ต้นไม้ลึก 3 ชั้น และ Logistic Regression เพื่อดูว่าความสามารถในการแยกต้องอาศัยแบบจำลองที่ซับซ้อนหรือไม่
    \item \textbf{การสลับป้ายกำกับ (label-shuffle)} ใช้แถวจริงชุดเดียวกันแต่สลับป้ายกำกับคอร์ปัสก่อนการแบ่งข้อมูล หากยังแยกได้แสดงว่าไปป์ไลน์รับรู้ป้ายกำกับผ่านช่องทางอื่นที่ไม่ใช่คุณลักษณะ
    \item \textbf{กลุ่มควบคุมว่าง (null control)} แบ่งคอร์ปัสเดียวออกเป็นสองครึ่งแบบสุ่ม เพื่อตรวจว่าวิธีการสร้างการแยกขึ้นเองหรือไม่
    \item \textbf{กลุ่มควบคุมภายในคอร์ปัส} แยกคนละไฟล์หรือคนละช่วงเวลาภายในคอร์ปัสเดียวกัน ด้วยแบบจำลองทั้ง 7 รูปแบบและจำนวนแถวต่อฝั่งเท่ากัน และกรณีที่ตรึงเครื่องมือโจมตีให้เป็นตัวเดียวกันแล้วเปลี่ยนเฉพาะคอร์ปัส
    \item \textbf{ระเบียนซ้ำ} นับระเบียนที่ซ้ำกันทุกค่าคร่อมชุดฝึกและชุดทดสอบ และระเบียนใกล้ซ้ำ (near-duplicate) ซึ่งนิยามเป็นแถวทดสอบที่อยู่ใกล้แถวฝึกมากกว่าระยะที่แถวฝึกอยู่ห่างกันเองถึง 99.9\% ของกรณี ระดับที่เกิดขึ้นโดยบังเอิญจึงเท่ากับ 0.10\%
\end{itemize}

\subsection{การวิเคราะห์สาเหตุของการทำนายพลาดด้วย SHAP}
\label{sec:cd-shap-method}

สมมติฐานที่ทดสอบคือ \emph{เมื่อนำแบบจำลองไปใช้ข้ามคอร์ปัส ระเบียนโจมตีที่ถูกทำนายผิดเป็น Benign ถูกผลักไปทาง Benign ผ่านกลุ่มคุณลักษณะที่บ่งชี้คอร์ปัสมากกว่าระเบียนที่ทำนายถูก} กลุ่มคุณลักษณะดังกล่าว (F) เลือกจากคุณลักษณะ 60 ตัว (77 คอลัมน์หลังรวมคอลัมน์ที่มีค่าซ้ำกัน) โดยใช้ผลการตรวจสอบเชิงปฏิปักษ์ และนับเฉพาะคุณลักษณะที่ผ่านทั้งเกณฑ์ permutation importance และเกณฑ์ SHAP ได้กลุ่มหลัก (consensus) 13 ตัว ได้แก่ min\_seg\_size\_forward, RST Flag Count, Init\_Win\_bytes\_forward, Destination Port, Max Packet Length, act\_data\_pkt\_fwd, Fwd IAT Min, Init\_Win\_bytes\_backward, Flow IAT Min, Bwd IAT Min, Min Packet Length, Bwd Packet Length Min และ Bwd Packet Length Std กลุ่มนี้ถูกกำหนดไว้ก่อนการทดสอบและไม่ถูกเปลี่ยนภายหลัง

การวิเคราะห์ใช้แบบจำลองต้นไม้ 4 รูปแบบ (LightGBM, Random Forest, XGBoost และ CatBoost) บนการแบ่งตามลำดับการจับข้อมูล 70/30 ค่าเริ่มต้นการสุ่ม 42--46 และใช้ไม่เกิน 1,000 แถวต่อคลาสต่อกลุ่มต่อรอบ สำหรับแต่ละแถวคำนวณ \textit{s} คือสัดส่วนของค่า SHAP ที่เป็นบวกต่อผลลัพธ์ Benign ซึ่งมาจากคุณลักษณะใน F การเปรียบเทียบในเกณฑ์ตัดสินสุดท้ายเป็นแบบ \textbf{ภายในคอร์ปัสปลายทาง} (within-target) คือ

\begin{equation}
    G_{\mathrm{within}} = \bar{s}(\mathrm{missed}) - \bar{s}(\mathrm{caught_{tgt}})
\end{equation}

โดยที่ missed คือระเบียนโจมตีในชุดทดสอบของคอร์ปัสปลายทางที่ถูกทำนายเป็น Benign และ caught\textsubscript{tgt} คือระเบียนโจมตีคลาสเดียวกันในคอร์ปัสเดียวกันที่ถูกทำนายถูก การเทียบภายในคอร์ปัสปลายทางตัดผลของการที่ข้อมูลมาจากต่างคอร์ปัสออก ซึ่งการเทียบกับระเบียนที่ทำนายถูกในคอร์ปัสต้นทาง (ที่ใช้ในรอบก่อนหน้า) แยกไม่ได้ ค่า \textit{G} ถูกเทียบกับกลุ่มคุณลักษณะสุ่มขนาดเท่ากัน 2,000 กลุ่ม (\textit{p} แบบด้านเดียว) กรณีที่นับได้ต้องมีระเบียน missed อย่างน้อย 100 แถว และ caught\textsubscript{tgt} อย่างน้อย 30 แถวในทุกรอบ และกรณีหนึ่งจะถือว่าผ่านเมื่อ \textit{G}\textsubscript{within} > 0 และ \textit{p} ต่ำกว่า 0.05 หารด้วยจำนวนกรณีที่นับได้ของแบบจำลองนั้น (Bonferroni) เกณฑ์ทั้งหมดถูกระบุไว้ก่อนการวิเคราะห์ (criteria specified before analysis) ในเอกสารภายในของโครงงานที่ระบุวันที่ไว้ก่อนการรัน แต่มิได้ลงทะเบียนกับหน่วยงานหรือระบบภายนอก จึงไม่ถือเป็นการลงทะเบียนล่วงหน้า (pre-registration) ในความหมายที่เป็นทางการ

\section{การควบคุมให้การทดลองทำซ้ำได้}

เพื่อให้ผลการทดลองสามารถทำซ้ำและตรวจสอบย้อนกลับได้ โครงงานนี้กำหนดมาตรการควบคุมดังนี้

\begin{enumerate}
    \item สำหรับการทดลองหลัก กำหนดค่าเริ่มต้นของการสุ่ม (random seed) เท่ากับ 42 เพื่อให้สามารถทำซ้ำผลได้ ส่วนการทดลองเสริมข้ามชุดข้อมูลใช้ค่าเริ่มต้นการสุ่ม 5 ค่า ได้แก่ 42--46 เพื่อวัดความแปรปรวนจากการสุ่มของแบบจำลองและการแบ่งข้อมูลแบบสุ่ม ตามที่อธิบายในหัวข้อ 3.8.4 สำหรับการแบ่งตามลำดับการจับข้อมูล ชุดฝึกและชุดทดสอบคงเดิมทั้ง 5 รอบ จึงวัดความแปรปรวนจากความสุ่มของแบบจำลอง ส่วนการสุ่มตัวอย่างคอร์ปัสละ 300,000 ระเบียนใช้ seed 42 เพียงครั้งเดียวและใช้ตัวอย่างชุดเดิมทุกรอบ
    \item ตรึงรุ่นของไลบรารีที่ใช้ทั้งหมดไว้ในไฟล์รายการที่ต้องติดตั้ง
    \item คำนวณค่าลายนิ้วมือดิจิทัลแบบ SHA-256 ของชุดข้อมูลที่ผ่านการทำความสะอาดแล้ว และบันทึกไว้พร้อมกับผลลัพธ์ของทุกการทดลอง เพื่อยืนยันว่าผลลัพธ์แต่ละชุดมาจากข้อมูลชุดเดียวกัน
    \item บันทึกไฟล์ประกาศการแบ่งข้อมูล (split manifest) ซึ่งระบุจำนวนตัวอย่างที่คาดหมายในแต่ละคลาส และมีขั้นตอนตรวจสอบที่สร้างการแบ่งข้อมูลขึ้นใหม่จากแคชแล้วเปรียบเทียบกับค่าที่ประกาศไว้
    \item บันทึกรายชื่อและลำดับของคุณลักษณะไว้พร้อมกับไฟล์แบบจำลองแต่ละไฟล์ เพื่อป้องกันความคลาดเคลื่อนของโครงสร้างข้อมูลระหว่างขั้นตอนการฝึกและขั้นตอนการทำนาย
    \item จัดทำชุดทดสอบซอฟต์แวร์แบบอัตโนมัติครอบคลุมไปป์ไลน์การประมวลผลข้อมูล ส่วนบริการเบื้องหลัง และวงจรชีวิตของไฟล์ผลลัพธ์
\end{enumerate}

    \chapter{ผลการทดลองและระบบต้นแบบ}
\label{chapter:result}

บทนี้รายงานผลการทดลองตามระเบียบวิธีในบทที่ \ref{chapter:experiment} แบ่งเป็นสามส่วน ได้แก่ ผลการทดลองหลักบนชุดข้อมูล CICIDS2017 แบบ 9 คลาส ระบบต้นแบบ และผลการทดลองเสริมเรื่องการส่งผ่านข้ามชุดข้อมูล ส่วนท้ายของบทรวบรวมข้อค้นพบสำคัญพร้อมตัวเลขไว้ในหน้าเดียว

%======================================================================
\section{ผลการทดลองหลักบนชุดข้อมูล CICIDS2017}
\label{sec:main-results}

\subsection{ประสิทธิภาพการจำแนกประเภท}
\label{sec:main-classification}

ตารางที่ \ref{tab:main-classification} แสดงผลของแบบจำลองทั้ง 7 รูปแบบบนชุดทดสอบ 749,401 ระเบียน ซึ่งแบ่งตามลำดับการจับข้อมูล (capture-order) แบบ 70/30 ภายในทุกกลุ่ม (ไฟล์, คลาส) คอลัมน์ Macro F1 (7 คลาส) คิดเฉพาะคลาสที่มีตัวอย่างในชุดทดสอบไม่น้อยกว่า 30 ระเบียน คือไม่รวม Heartbleed และ Infiltration

ค่า Macro F1-score ที่ใช้เป็นผลหลักและเกณฑ์จัดอันดับในงานจำแนก 9 คลาส คือค่าเฉลี่ย F1 ของทั้ง 9 คลาสโดยให้น้ำหนักแต่ละคลาสเท่ากัน ดังนั้นค่า 0.8873 ของ CatBoost เป็นค่าของทั้ง 9 คลาส ส่วนค่า 0.8613 ในคอลัมน์ 7 คลาสเป็นตัวชี้วัดประกอบเพื่อตรวจความไวของคะแนนต่อคลาสที่มีตัวอย่างน้อย โดยตัดเฉพาะ Heartbleed (ชุดทดสอบ 4 ระเบียน) และ Infiltration (11 ระเบียน) ออกจากการเฉลี่ย ไม่ได้ฝึกแบบจำลองใหม่หรือลบระเบียนเหล่านี้ออกจากชุดทดสอบ ผลรายคลาสยังรายงานครบทั้ง 9 คลาส และคะแนนจากคลาสที่มีตัวอย่างน้อยต้องตีความอย่างระมัดระวัง

ตัวชี้วัด 7 คลาสข้างต้นแตกต่างจาก Macro F1 ของการทดลองเสริมข้ามชุดข้อมูลในหัวข้อ \ref{sec:crossdataset} ซึ่งฝึกและประเมินบนโครงร่างคลาสร่วม 7 คลาส ได้แก่ Benign, DDoS, DoS, Brute Force, Bot, Infiltration และ Web Attack โดยตัด PortScan และ Heartbleed ที่ไม่มีคลาสคู่กันใน CSE-CIC-IDS2018 ออก จึงไม่ควรเปรียบเทียบค่าของทั้งสองกรณีโดยตรง แม้มีจำนวนคลาสเท่ากัน

% แหล่งข้อมูล: results/cicids2017_temporal_v1/report.md
\begin{table}[ht]
    \caption{ประสิทธิภาพการจำแนกประเภทของแบบจำลองทั้ง 7 รูปแบบบนชุดทดสอบของ CICIDS2017 (9 คลาส)}
    \label{tab:main-classification}
    \small
    \begin{tabularx}{\linewidth}{@{}l *{6}{>{\raggedleft\arraybackslash}X}@{}}
        \hline
        \textbf{แบบจำลอง} & \textbf{Accuracy} & \textbf{Balanced Acc.} & \textbf{Macro F1 (9 คลาส)} & \textbf{Macro F1 (7 คลาส)} & \textbf{Weighted F1} & \textbf{MCC} \\
        \hline
        CatBoost & 0.9953 & 0.8970 & \textbf{0.8873} & 0.8613 & 0.9952 & 0.9844 \\
        LightGBM & 0.9903 & 0.8576 & 0.8774 & 0.8627 & 0.9898 & 0.9676 \\
        XGBoost & 0.9906 & 0.8889 & 0.8773 & 0.8542 & 0.9902 & 0.9687 \\
        Random Forest & 0.9820 & 0.8165 & 0.8404 & 0.8378 & 0.9832 & 0.9397 \\
        Stacking Ensemble & 0.9581 & 0.8852 & 0.6822 & 0.7092 & 0.9712 & 0.8757 \\
        Logistic Regression & 0.8428 & 0.8338 & 0.5574 & 0.6348 & 0.9030 & 0.6651 \\
        Multi-Layer Perceptron & 0.9331 & 0.7112 & 0.5369 & 0.6767 & 0.9443 & 0.7895 \\
        \hline
    \end{tabularx}
\end{table}

CatBoost ได้ค่า Macro F1-score สูงสุดที่ 0.8873 ตามด้วย LightGBM (0.8774) และ XGBoost (0.8773) ซึ่งต่างกันไม่ถึง 0.01 แบบจำลองทุกตัวมีค่า Accuracy สูงกว่าแบบจำลองฐานที่ทำนาย BENIGN ทุกกรณี (0.8296) แต่ Logistic Regression สูงกว่าเพียงเล็กน้อย (0.8428) ซึ่งยืนยันว่าค่า Accuracy เพียงค่าเดียวไม่เพียงพอต่อการตัดสินบนข้อมูลที่ไม่สมดุล

รูปที่ \ref{fig:main-cm-catboost} แสดงผลการจำแนกของ CatBoost บนชุดทดสอบของรอบหลัก แถวแสดงคลาสจริงและคอลัมน์แสดงคลาสที่แบบจำลองทำนาย จาก Bot ทั้ง 585 ระเบียน แบบจำลองทำนายถูก 94 ระเบียน (ความไว 0.1607) และทำนายอีก 491 ระเบียนเป็น BENIGN ในอีกทิศทางหนึ่ง มี BENIGN 233 ระเบียนที่ถูกทำนายเป็น Bot สำหรับ DoS แบบจำลองทำนายผิดเป็น BENIGN 1,431 ระเบียน เป็น PortScan 913 ระเบียน และเป็น Web Attack 151 ระเบียน เมทริกซ์ของอีก 6 แบบจำลองแสดงในภาคผนวก \ref{appendix:main-confusion-matrices}

% แหล่งข้อมูล: results/cicids2017_temporal_v1/catboost_confusion_matrix.png
\begin{figure}[!htbp]
    \centering
    \includegraphics[width=0.92\linewidth]{catboost_confusion_matrix.png}
    \caption{เมทริกซ์ความสับสนของ CatBoost จากการทดลองหลักบนชุดทดสอบ CICIDS2017}
    \label{fig:main-cm-catboost}
\end{figure}
\FloatBarrier

\subsection{ความสามารถในการตรวจจับและผลกระทบเชิงปฏิบัติการ}

ตารางที่ \ref{tab:main-detection} ยุบผลการทำนายเป็นกรณีทวิภาคระหว่างการจราจรปกติและการโจมตี CatBoost ตรวจจับการโจมตีได้ 98.40\% โดยมีการแจ้งเตือนผิดพลาด 404 ครั้ง หรือ 6.5 ครั้งต่อการจราจรปกติหนึ่งหมื่นระเบียน ขณะที่ Logistic Regression มีอัตราการตรวจจับใกล้เคียงกัน (97.46\%) แต่แจ้งเตือนผิดพลาดถึง 105,258 ครั้ง

\begin{table}[ht]
    \caption{ความสามารถในการตรวจจับการโจมตีแบบทวิภาคบนชุดทดสอบของ CICIDS2017}
    \label{tab:main-detection}
    \small
    \begin{tabularx}{\linewidth}{@{}l *{5}{>{\raggedleft\arraybackslash}X}@{}}
        \hline
        \textbf{แบบจำลอง} & \textbf{อัตราตรวจจับ} & \textbf{FPR} & \textbf{แจ้งเตือนผิด} & \textbf{ต่อปกติ 10,000} & \textbf{การโจมตีที่พลาด} \\
        \hline
        CatBoost & 0.9840 & 0.00065 & 404 & 6.5 & 2,049 \\
        Logistic Regression & 0.9746 & 0.16931 & 105,258 & 1,693.1 & 3,244 \\
        Stacking Ensemble & 0.9661 & 0.04072 & 25,317 & 407.2 & 4,329 \\
        XGBoost & 0.9571 & 0.00038 & 236 & 3.8 & 5,483 \\
        LightGBM & 0.9513 & 0.00022 & 136 & 2.2 & 6,220 \\
        Random Forest & 0.9277 & 0.00530 & 3,298 & 53.0 & 9,232 \\
        Multi-Layer Perceptron & 0.8459 & 0.04721 & 29,347 & 472.1 & 19,677 \\
        \hline
    \end{tabularx}
\end{table}

เมื่อพิจารณารายคลาส คลาส Bot เป็นคลาสที่ตรวจจับได้แย่ที่สุดในทุกแบบจำลอง โดยมีค่าความไวเพียง 0.15--0.16 ส่วนคลาส DDoS และ PortScan มีค่าความไวเกิน 0.99 และ Brute Force เกิน 0.95 ในทุกแบบจำลอง คลาส Heartbleed (4 ระเบียน) และ Infiltration (11 ระเบียน) มีตัวอย่างในชุดทดสอบน้อยเกินกว่าจะตีความค่าความไวเป็นค่าที่เสถียรได้

\subsection{ประสิทธิภาพเชิงคำนวณและการจัดอันดับ}

ตารางที่ \ref{tab:main-efficiency} แสดงเวลาแฝงของการทำนายบนหน่วยประมวลผลกลาง โดยวัดเป็นชุดละ 1,000 โฟลว์ CatBoost มีเวลาแฝงที่ควอนไทล์ที่ 95 เท่ากับ 8.39 มิลลิวินาที ต่ำเป็นอันดับสองรองจาก Logistic Regression และมีขนาดไฟล์แบบจำลองเพียง 2.5 เมกะไบต์

\begin{table}[ht]
    \caption{ประสิทธิภาพเชิงคำนวณของแบบจำลอง (การทำนายบนหน่วยประมวลผลกลาง ชุดละ 1,000 โฟลว์)}
    \label{tab:main-efficiency}
    \small
    \begin{tabularx}{\linewidth}{@{}l *{5}{>{\raggedleft\arraybackslash}X}@{}}
        \hline
        \textbf{แบบจำลอง} & \textbf{เวลาฝึก (วินาที)} & \textbf{p50 (มิลลิวินาที)} & \textbf{p95 (มิลลิวินาที)} & \textbf{โฟลว์ต่อวินาที} & \textbf{ขนาด (MB)} \\
        \hline
        Logistic Regression & 653.7 & 4.30 & 5.15 & 232,391 & 0.0 \\
        CatBoost & 44.6 & 7.58 & 8.39 & 131,975 & 2.5 \\
        Multi-Layer Perceptron & 818.7 & 10.05 & 12.50 & 99,545 & 0.7 \\
        XGBoost & 154.6 & 43.42 & 54.10 & 23,030 & 15.8 \\
        Random Forest & 314.6 & 85.94 & 104.12 & 11,636 & 124.5 \\
        Stacking Ensemble & 2,754.7 & 119.79 & 137.95 & 8,348 & 109.3 \\
        LightGBM & 549.6 & 111.40 & 144.02 & 8,977 & 54.9 \\
        \hline
    \end{tabularx}
\end{table}

เมื่อนำเกณฑ์การจัดอันดับสามชุดที่ระบุไว้ก่อนการวิเคราะห์ผลมาใช้ CatBoost เป็นอันดับหนึ่งในทั้งสามเกณฑ์ ได้แก่ อันดับภาพรวม (Macro F1-score 0.8873) อันดับด้านความมั่นคงปลอดภัย (ความไวเฉลี่ยแบบ macro 0.8970 ภายใต้ FPR แบบทวิภาคไม่เกิน 0.01) และอันดับด้านการนำไปใช้งาน (p95 8.39 มิลลิวินาที ในกลุ่มแบบจำลองที่ Macro F1-score ต่างจากอันดับหนึ่งไม่เกิน 0.02) CatBoost จึงถูกเลือกเป็นแบบจำลองแชมเปียนของโครงงาน

\subsection{ผลการอธิบาย CatBoost ที่ฝึกซ้ำด้วย SHAP}

การวิเคราะห์ด้วย SHAP ตามวิธีในหัวข้อ \ref{sec:shap} ใช้ CatBoost ที่ฝึกซ้ำด้วยค่าไฮเปอร์พารามิเตอร์ในตารางที่ \ref{tab:best-params} บนชุดฝึกเดิมทั้ง 1,748,579 ระเบียน แบบจำลองที่ฝึกซ้ำได้ค่า Macro F1-score 0.8858 และ FPR 0.00063 บนชุดทดสอบเดิม ส่วน CatBoost ของรอบหลักที่ใช้สร้างเมทริกซ์ความสับสนในรูปที่ \ref{fig:main-cm-catboost} ได้ค่า Macro F1-score 0.8873 และ FPR 0.00065 ตามตารางที่ \ref{tab:main-classification} และ \ref{tab:main-detection} ค่า SHAP คำนวณด้วย TreeSHAP บนระเบียนจากชุดทดสอบที่สุ่มแยกตามคลาส คลาสละไม่เกิน 200 ระเบียน (Heartbleed มี 4 ระเบียนและ Infiltration มี 11 ระเบียนตามจำนวนที่มีในชุดทดสอบ) เพื่อไม่ให้ผลถูกครอบงำด้วยคลาส Benign ซึ่งมีสัดส่วนมากกว่า 80\% ของชุดทดสอบ

\begin{figure}[ht]
    \centering
    \includegraphics[width=0.8\linewidth]{fig_shap_main_bar.png}
    \caption{คุณลักษณะ 20 อันดับแรกของแบบจำลอง CatBoost เรียงตามค่าเฉลี่ยของค่าสัมบูรณ์ของค่า SHAP รวมทุกคลาส}
    \label{fig:shap-main-bar}
\end{figure}

รูปที่ \ref{fig:shap-main-bar} แสดงว่าในภาพรวม คุณลักษณะที่มีอิทธิพลสูงสุดคือ Destination Port (17.9\% ของผลรวมค่า SHAP) ตามด้วยขนาดหน้าต่าง TCP เริ่มต้นของทั้งสองทิศ คือ Init\_Win\_bytes\_forward (9.8\%) และ Init\_Win\_bytes\_backward (9.5\%) จากนั้นจึงเป็น min\_seg\_size\_forward (5.0\%) Flow Bytes/s (4.6\%) และกลุ่มช่วงเวลาระหว่างแพ็กเก็ตต่ำสุด ได้แก่ Flow IAT Min (4.3\%) และ Fwd IAT Min (3.5\%) ตารางที่ \ref{tab:shap-main-class} แสดงคุณลักษณะ 3 อันดับแรกของแต่ละคลาส

\begin{table}[ht]
    \caption{คุณลักษณะ 3 อันดับแรกของแต่ละคลาสจากการวิเคราะห์ SHAP ของแบบจำลอง CatBoost (ร้อยละของผลรวมค่า SHAP ของคลาสนั้น)}
    \label{tab:shap-main-class}
    \footnotesize
    \begin{tabularx}{\linewidth}{@{}l X X X@{}}
        \hline
        \textbf{คลาส} & \textbf{อันดับ 1} & \textbf{อันดับ 2} & \textbf{อันดับ 3} \\
        \hline
        Benign & Destination Port (16.3) & Init\_Win\_bytes\_forward (14.8) & Init\_Win\_bytes\_backward (9.6) \\
        DoS & Init\_Win\_bytes\_forward (18.3) & Destination Port (17.5) & Init\_Win\_bytes\_backward (6.4) \\
        DDoS & Init\_Win\_bytes\_backward (12.0) & Init\_Win\_bytes\_forward (10.6) & min\_seg\_size\_forward (10.3) \\
        PortScan & Init\_Win\_bytes\_forward (15.1) & Init\_Win\_bytes\_backward (10.7) & Destination Port (10.2) \\
        Brute Force & Destination Port (44.9) & Init\_Win\_bytes\_forward (6.2) & Fwd IAT Min (5.5) \\
        Bot & Destination Port (42.5) & Flow IAT Min (5.9) & Init\_Win\_bytes\_backward (5.8) \\
        Web Attack & Total Length of Fwd Packets (11.3) & Fwd Packet Length Max (10.8) & Fwd IAT Mean (6.3) \\
        Infiltration\textsuperscript{*} & Destination Port (26.1) & Init\_Win\_bytes\_forward (9.2) & min\_seg\_size\_forward (6.7) \\
        Heartbleed\textsuperscript{*} & Destination Port (25.9) & Init\_Win\_bytes\_backward (10.0) & Init\_Win\_bytes\_forward (6.8) \\
        \hline
        \multicolumn{4}{@{}l}{\scriptsize \textsuperscript{*} คำนวณจากระเบียนทดสอบเพียง 11 และ 4 ระเบียนตามลำดับ จึงใช้ดูแนวโน้มเท่านั้น} \\
    \end{tabularx}
\end{table}

ในระดับรายตัวอย่าง รูปที่ \ref{fig:shap-main-dos} แสดงค่า SHAP ของแต่ละระเบียนต่อผลลัพธ์คลาส DoS แต่ละจุดคือหนึ่งระเบียน ตำแหน่งตามแนวนอนคือขนาดและทิศทางที่คุณลักษณะนั้นผลักผลการทำนายเข้าหาหรือออกจาก DoS ส่วนสีแสดงค่าของคุณลักษณะ (น้ำเงินคือค่าต่ำ แดงคือค่าสูง) ระเบียนที่มีค่า Init\_Win\_bytes\_forward ต่ำถูกผลักไปทาง DoS อย่างชัดเจน (ค่า SHAP ประมาณ +2.5 ถึง +6) ขณะที่ระเบียนที่มีค่าสูงถูกผลักออก (ประมาณ −3) ส่วนค่า Destination Port ที่ต่ำ เช่น พอร์ต 80 ของเว็บเซิร์ฟเวอร์ ก็ผลักไปทาง DoS เช่นกัน

\begin{figure}[ht]
    \centering
    \includegraphics[width=0.72\linewidth]{fig_shap_main_dos.png}
    \caption{ค่า SHAP รายระเบียนของแบบจำลอง CatBoost ต่อผลลัพธ์คลาส DoS (15 คุณลักษณะแรก)}
    \label{fig:shap-main-dos}
\end{figure}

ข้อสังเกตจากผลการวิเคราะห์มีดังนี้

\begin{itemize}
    \item \textbf{แบบจำลองพึ่งพาหมายเลขพอร์ตปลายทางมาก} โดยเฉพาะ Brute Force และ Bot ที่ Destination Port คิดเป็นสัดส่วนเกือบครึ่งหนึ่งของค่า SHAP ทั้งหมด ซึ่งสอดคล้องกับลักษณะของการโจมตีที่มุ่งไปยังบริการเฉพาะ เช่น FTP (พอร์ต 21) และ SSH (พอร์ต 22) สำหรับ Brute Force และพอร์ต 8080 ซึ่ง Bot ใน CICIDS2017 มากกว่าครึ่งหนึ่งใช้ติดต่อกับเครื่องสั่งการ
    \item \textbf{คุณลักษณะระดับโพรโทคอล TCP มีบทบาทสูงในคลาสที่มีปริมาณมาก} ขนาดหน้าต่าง TCP เริ่มต้นและขนาดเซกเมนต์ต่ำสุดเป็นคุณลักษณะหลักของ DoS, DDoS และ PortScan ค่าเหล่านี้ขึ้นกับระบบปฏิบัติการและเครื่องมือที่ใช้สร้างทราฟฟิก มากกว่าพฤติกรรมของการโจมตีโดยตรง
    \item \textbf{Web Attack เป็นคลาสเดียวที่อาศัยคุณลักษณะด้านเนื้อหาและจังหวะเวลา} ได้แก่ ความยาวของข้อมูลขาไปและช่วงเวลาระหว่างแพ็กเก็ต ซึ่งสมเหตุสมผล เพราะการโจมตีเว็บใช้พอร์ต 80 เช่นเดียวกับทราฟฟิกปกติ
    \item \textbf{คุณลักษณะอันดับต้นซ้อนทับกับกลุ่มที่บ่งชี้ชุดข้อมูล} คุณลักษณะ 10 อันดับแรกในภาพรวมมี 8 ตัวที่อยู่ในกลุ่มคุณลักษณะที่บ่งชี้คอร์ปัส 13 ตัวของการทดลองเสริม (หัวข้อ \ref{sec:cd-shap-method}) แม้ทั้งสองการวิเคราะห์จะใช้โครงร่างคลาสและแบบจำลองต่างกัน ผลนี้ชี้ว่าคุณลักษณะที่แบบจำลองใช้จำแนกการโจมตีได้ดีภายใน CICIDS2017 เป็นกลุ่มเดียวกับคุณลักษณะที่มีค่าต่างกันมากระหว่างชุดข้อมูล ซึ่งเป็นคำอธิบายหนึ่งของผลการส่งผ่านที่จำกัดในหัวข้อ \ref{sec:cd-results}
\end{itemize}

%======================================================================
\FloatBarrier
\section{ระบบต้นแบบ}
\label{sec:prototype}

ระบบต้นแบบประกอบด้วยส่วนบริการเบื้องหลังที่พัฒนาด้วย FastAPI และส่วนติดต่อผู้ใช้ที่พัฒนาด้วย Next.js ส่วนติดต่อผู้ใช้แบ่งเป็นหน้าหลักดังนี้

\begin{itemize}
    \item \textbf{ภาพรวม} สรุปแบบจำลองแชมเปียนและตัวชี้วัดหลักของชุดผลลัพธ์ที่เลือก
    \item \textbf{ชุดข้อมูลและการสำรวจข้อมูล} แสดงการกระจายตัวของคลาสและผลการสำรวจข้อมูลเบื้องต้น
    \item \textbf{ประสิทธิภาพและการเปรียบเทียบ} แสดงตัวชี้วัดทั้งสี่มิติ เมทริกซ์ความสับสน และการจัดอันดับตามเกณฑ์ทั้งสามชุด
    \item \textbf{การอธิบายแบบจำลอง} แสดงคุณลักษณะที่มีอิทธิพลสูงสุดจากการวิเคราะห์ SHAP ทั้งภาพรวมและรายคลาส พร้อมค่าความเสถียรระหว่างรอบการสุ่ม
    \item \textbf{การทำนาย} รับไฟล์ข้อมูลโฟลว์ใหม่ ตรวจสอบโครงสร้างคอลัมน์เทียบกับรายชื่อคุณลักษณะที่บันทึกไว้พร้อมแบบจำลอง แล้วทำนายผลเป็นชุด
\end{itemize}

ภาพหน้าจอของแต่ละหน้าแสดงในรูปที่ \ref{fig:dash-overview} ถึง \ref{fig:dash-predict} ทุกภาพแสดงชุดผลลัพธ์ \texttt{cicids2017\_temporal\_v1} ของการทดลองหลัก (ชื่อโฟลเดอร์ในโค้ด ซึ่งใช้การแบ่งตามลำดับการจับข้อมูลตามหัวข้อ \ref{sec:split})

\begin{figure}[ht]
    \centering
    \includegraphics[width=\linewidth]{fig_dash_overview.png}
    \caption{หน้าภาพรวมของเว็บแดชบอร์ด แสดงแบบจำลองแชมเปียน ตัวชี้วัดหลัก และรายชื่อแบบจำลองในชุดผลลัพธ์}
    \label{fig:dash-overview}
\end{figure}

\begin{figure}[ht]
    \centering
    \includegraphics[width=\linewidth]{fig_dash_dataset.png}
    \caption{หน้าชุดข้อมูล แสดงจำนวนระเบียนและการกระจายตัวของคลาสในชุดฝึกและชุดทดสอบ}
    \label{fig:dash-dataset}
\end{figure}

\begin{figure}[ht]
    \centering
    \includegraphics[width=\linewidth]{fig_dash_performance.png}
    \caption{หน้าประสิทธิภาพของแบบจำลอง แสดงตัวชี้วัดรวมและค่าตัวชี้วัดรายคลาสของแบบจำลองที่เลือก}
    \label{fig:dash-performance}
\end{figure}

\begin{figure}[ht]
    \centering
    \includegraphics[width=\linewidth]{fig_dash_compare.png}
    \caption{หน้าเปรียบเทียบ แสดงตัวชี้วัดของแบบจำลองทั้ง 7 รูปแบบในตารางเดียว}
    \label{fig:dash-compare}
\end{figure}

\begin{figure}[ht]
    \centering
    \includegraphics[width=\linewidth]{fig_dash_predict.png}
    \caption{หน้าการทำนาย สำหรับเลือกแบบจำลองและอัปโหลดไฟล์ CSV ของโฟลว์ใหม่เพื่อทำนายผลเป็นชุด}
    \label{fig:dash-predict}
\end{figure}

%======================================================================
\FloatBarrier
\section{ผลการทดลองเสริม: การส่งผ่านข้ามชุดข้อมูล}
\label{sec:cd-results}

ผลในหัวข้อนี้ทั้งหมดมาจาก protocol\_v3 (ชุดฝึก 70\% / ชุดทดสอบ 30\%) ตามที่อธิบายในหัวข้อ \ref{sec:crossdataset} ผลของ protocol\_v2 (ชุดฝึก 60\%) ถูกอ้างเฉพาะในการวิเคราะห์ความไวต่อสัดส่วนการแบ่ง (หัวข้อ \ref{sec:cd-ssh})

\subsection{ผลหลัก}

ตารางที่ \ref{tab:cd-main-random} และ \ref{tab:cd-main-capture} แสดงค่า Macro F1-score ของการส่งผ่าน (ค่าเฉลี่ย ± ส่วนเบี่ยงเบนมาตรฐานจาก 5 รอบ) ค่า Recovery และค่าเพดานของแต่ละแบบจำลอง ค่าเพดานในแต่ละทิศคือค่าที่แบบจำลองชนิดเดียวกันทำได้ภายในคอร์ปัสปลายทาง ค่าพื้นของทิศ 2017\dirarrow{}2018 คือ 0.1337 และของทิศ 2018\dirarrow{}2017 คือ 0.1294 ตัวหนาคือค่า Macro F1-score สูงสุดของแต่ละทิศ

% แหล่งข้อมูล: docs/experiments/crossdataset/round6_split_70_30/data/summary_macro_f1.csv และ summary_gap.csv
\begin{table}[ht]
    \caption{ผลการส่งผ่านข้ามชุดข้อมูลภายใต้การแบ่งแบบสุ่ม (protocol\_v3)}
    \label{tab:cd-main-random}
    \footnotesize
    \begin{tabularx}{\linewidth}{@{}l c *{2}{>{\raggedleft\arraybackslash}X} c *{2}{>{\raggedleft\arraybackslash}X}@{}}
        \hline
        & \multicolumn{3}{c}{\textbf{2017\dirarrow{}2018 (ค่าพื้น 0.1337)}} & \multicolumn{3}{c}{\textbf{2018\dirarrow{}2017 (ค่าพื้น 0.1294)}} \\
        \textbf{แบบจำลอง} & \textbf{Macro F1} & \textbf{Rec.\ (\%)} & \textbf{เพดาน} & \textbf{Macro F1} & \textbf{Rec.\ (\%)} & \textbf{เพดาน} \\
        \hline
        LightGBM & \textbf{0.2850} ± 0.0181 & 20.6 & 0.8696 & 0.2860 ± 0.0524 & 18.4 & 0.9795 \\
        XGBoost & 0.2727 ± 0.0633 & 18.9 & 0.8676 & 0.2446 ± 0.0179 & 13.7 & 0.9715 \\
        CatBoost & 0.1692 ± 0.0073 & 4.9 & 0.8547 & 0.2195 ± 0.0176 & 10.7 & 0.9686 \\
        Random Forest & 0.1743 ± 0.0086 & 5.6 & 0.8625 & 0.1492 ± 0.0024 & 2.4 & 0.9612 \\
        Stacking & 0.2836 ± 0.0094 & 21.3 & 0.8383 & \textbf{0.3187} ± 0.0296 & 22.4 & 0.9727 \\
        MLP & 0.1489 ± 0.0068 & 2.2 & 0.8326 & 0.2256 ± 0.0173 & 13.0 & 0.8679 \\
        Logistic Regression & 0.1381 ± 0.0007 & 0.7 & 0.7811 & 0.2556 ± 0.0271 & 25.5 & 0.6243 \\
        \hline
        \textbf{เฉลี่ย 7 แบบจำลอง} & & 10.6 & & & 15.2 & \\
        \hline
    \end{tabularx}
\end{table}

\begin{table}[ht]
    \caption{ผลการส่งผ่านข้ามชุดข้อมูลภายใต้การแบ่งตามลำดับการจับข้อมูล (protocol\_v3)}
    \label{tab:cd-main-capture}
    \footnotesize
    \begin{tabularx}{\linewidth}{@{}l c *{2}{>{\raggedleft\arraybackslash}X} c *{2}{>{\raggedleft\arraybackslash}X}@{}}
        \hline
        & \multicolumn{3}{c}{\textbf{2017\dirarrow{}2018 (ค่าพื้น 0.1337)}} & \multicolumn{3}{c}{\textbf{2018\dirarrow{}2017 (ค่าพื้น 0.1294)}} \\
        \textbf{แบบจำลอง} & \textbf{Macro F1} & \textbf{Rec.\ (\%)} & \textbf{เพดาน} & \textbf{Macro F1} & \textbf{Rec.\ (\%)} & \textbf{เพดาน} \\
        \hline
        LightGBM & 0.3132 ± 0.0569 & 24.5 & 0.8675 & \textbf{0.4435} ± 0.0041 & 42.0 & 0.8767 \\
        XGBoost & \textbf{0.3368} ± 0.0599 & 28.0 & 0.8589 & 0.3524 ± 0.0187 & 30.4 & 0.8617 \\
        CatBoost & 0.1735 ± 0.0017 & 5.5 & 0.8516 & 0.2176 ± 0.0042 & 11.9 & 0.8711 \\
        Random Forest & 0.1717 ± 0.0001 & 5.3 & 0.8466 & 0.1631 ± 0.0056 & 4.6 & 0.8550 \\
        Stacking & 0.3322 ± 0.0600 & 28.6 & 0.8272 & 0.3567 ± 0.0217 & 30.9 & 0.8648 \\
        MLP & 0.1753 ± 0.0015 & 6.1 & 0.8145 & 0.3005 ± 0.0219 & 27.5 & 0.7511 \\
        Logistic Regression & 0.1496 ± 0.0000 & 2.6 & 0.7534 & 0.3506 ± 0.0000 & 50.1 & 0.5706 \\
        \hline
        \textbf{เฉลี่ย 7 แบบจำลอง} & & 14.4 & & & 28.2 & \\
        \hline
    \end{tabularx}
\end{table}

\begin{figure}[ht]
    \centering
    \includegraphics[width=\linewidth]{fig_cd_split_modes.png}
    \caption{ค่า Macro F1-score ของการส่งผ่านภายใต้การแบ่งแบบสุ่มและการแบ่งตามลำดับการจับข้อมูล (เส้นประคือค่าพื้นของชุดทดสอบปลายทาง)}
    \label{fig:cd-split-modes}
\end{figure}

ข้อสังเกตจากตารางที่ \ref{tab:cd-main-random}, \ref{tab:cd-main-capture} และรูปที่ \ref{fig:cd-split-modes} มีดังนี้

\begin{itemize}
    \item \textbf{แบบจำลองที่ได้ค่าสูงสุดในแต่ละทิศ} ภายใต้การแบ่งตามลำดับการจับข้อมูล ทิศ 2017\dirarrow{}2018 คือ XGBoost (0.3368) ทิศ 2018\dirarrow{}2017 คือ LightGBM (0.4435) ภายใต้การแบ่งแบบสุ่ม ทิศ 2017\dirarrow{}2018 คือ LightGBM (0.2850) ทิศ 2018\dirarrow{}2017 คือ Stacking (0.3187) ค่าสูงสุดทุกกรณียังห่างจากค่าเพดาน (0.57--0.98) มาก
    \item \textbf{การแบ่งแบบสุ่มกับการแบ่งตามลำดับการจับข้อมูล} ค่าเพดานของ CICIDS2017 ภายใต้การแบ่งแบบสุ่มสูงกว่าการแบ่งตามลำดับการจับข้อมูลเฉลี่ย 0.099 (0.9065 เทียบกับ 0.8073) ขณะที่ของ CSE-CIC-IDS2018 ต่างกันเพียง 0.012 (0.8438 เทียบกับ 0.8314) ผลนี้แสดงความแตกต่างของค่าเพดานระหว่างสองวิธีแบ่งข้อมูล โดยยังไม่สามารถตัดสินว่าค่าใดสะท้อนประสิทธิภาพการใช้งานจริงได้ดีกว่า เนื่องจากการแบ่งตามลำดับการจับข้อมูลไม่ใช่การกันข้อมูลตามเวลาประทับ และองค์ประกอบของเครื่องมือโจมตีย่อยมีผลต่อคะแนนด้วย ค่า Recovery จึงต้องอ่านคู่กับค่าดิบและค่าเพดานของแต่ละวิธีแบ่งข้อมูล
    \item \textbf{ความไม่สมมาตร} ภายใต้การแบ่งตามลำดับการจับข้อมูล ค่า Recovery เฉลี่ยของทิศ 2018\dirarrow{}2017 (28.2\%) สูงกว่าทิศ 2017\dirarrow{}2018 (14.4\%) อยู่ 13.8 จุด และทิศ 2018\dirarrow{}2017 ดีกว่าใน 6 จาก 7 แบบจำลอง (Random Forest ต่างกันไม่ถึง 1 จุด) ภายใต้การแบ่งแบบสุ่ม ค่า Recovery เฉลี่ยของทิศ 2018\dirarrow{}2017 และ 2017\dirarrow{}2018 เท่ากับ 15.2\% และ 10.6\% ตามลำดับ ความไม่สมมาตรที่สังเกตได้จึงมีขนาดเล็กลงเป็น 4.6 จุดเปอร์เซ็นต์ และทิศ 2018\dirarrow{}2017 ดีกว่าใน 4 จาก 7 แบบจำลอง ความไม่สมมาตรจึงยังคงอยู่ แต่ขนาดขึ้นกับวิธีแบ่งข้อมูลและสัดส่วนการแบ่ง (หัวข้อ \ref{sec:cd-ssh})
    \item \textbf{การจัดอันดับต้องดูค่าดิบคู่กับค่า Recovery} Logistic Regression ได้ค่า Recovery สูงสุดในทิศ 2018\dirarrow{}2017 (50.1\%) แต่ค่า Macro F1-score ดิบ (0.3506) ต่ำกว่า LightGBM (0.4435) ที่ได้ Recovery 42.0\% เพราะค่าเพดานของ Logistic Regression บน CICIDS2017 ต่ำเพียง 0.5706 การใช้ค่า Recovery เพียงค่าเดียวจะจัดอันดับผิด
    \item \textbf{ระดับของค่าดิบ} ในทิศ 2017\dirarrow{}2018 มี 4 จาก 7 แบบจำลอง (CatBoost, Random Forest, MLP และ Logistic Regression) ที่ได้ค่า Macro F1-score เพียง 0.15--0.18 ใกล้กับค่าพื้น 0.1337
\end{itemize}

\subsection{ผลรายคลาส}
\label{sec:cd-per-class}

ตารางที่ \ref{tab:cd-per-class} และรูปที่ \ref{fig:cd-per-class} แสดงค่า F1 รายคลาสภายใต้การแบ่งตามลำดับการจับข้อมูล ค่าเพดานในตารางเป็นค่าเฉลี่ยของ 7 แบบจำลอง ส่วนในรูปเป็นค่าสูงสุดของ 7 แบบจำลอง

% แหล่งข้อมูล: round6_split_70_30/data/summary_per_class.csv และ report/protocol_v3_per_class.md
\begin{table}[ht]
    \caption{ค่า F1 รายคลาสของการส่งผ่านภายใต้การแบ่งตามลำดับการจับข้อมูล (protocol\_v3)}
    \label{tab:cd-per-class}
    \footnotesize
    \begin{tabularx}{\linewidth}{@{}l l l c c X@{}}
        \hline
        \textbf{คลาส} & \textbf{สูงสุด 2017\dirarrow{}2018} & \textbf{สูงสุด 2018\dirarrow{}2017} & \textbf{เพดาน 2017} & \textbf{เพดาน 2018} & \textbf{การส่งผ่าน} \\
        \hline
        Brute Force & 0.929 (Stacking) & 0.969 (XGBoost) & 0.990 & 0.999 & ดีในบางแบบจำลอง \\
        DoS & 0.279 (MLP) & 0.590 (CatBoost) & 0.910 & 0.939 & บางส่วน \\
        DDoS & 0.011 (MLP) & 0.518 (LR) & 0.993 & 0.991 & แย่ \\
        Bot & 0.358 (XGBoost) & 0.869 (LightGBM) & 0.208 & 0.998 & ขึ้นกับแบบจำลอง \\
        Web Attack & 0.002 (XGBoost) & 0.018 (LightGBM) & 0.803 & 0.782 & แย่ \\
        Benign & 0.942 (XGBoost) & 0.954 (LR) & 0.987 & 0.991 & ดี \\
        Infiltration & 0.000 & 0.000 & 0.760 & 0.120 & รายงานไม่ได้ \\
        \hline
    \end{tabularx}
\end{table}

\begin{figure}[ht]
    \centering
    \includegraphics[width=0.72\linewidth]{fig_cd_per_class.png}
    \caption{ค่า F1 สูงสุดจาก 7 แบบจำลองแยกตามคลาส ภายใต้การแบ่งตามลำดับการจับข้อมูล (\textsuperscript{*} Infiltration มีตัวอย่างทดสอบใน CICIDS2017 เพียง 11 ระเบียน)}
    \label{fig:cd-per-class}
\end{figure}

\begin{itemize}
    \item \textbf{Brute Force} ส่งผ่านได้ดีเมื่อแบบจำลองเป็น LightGBM, XGBoost หรือ Stacking ในทิศ 2017\dirarrow{}2018 (0.77--0.93) และในทิศ 2018\dirarrow{}2017 มี 5 แบบจำลองที่ได้ตั้งแต่ 0.5 ขึ้นไป แต่ CatBoost และ Random Forest ได้ไม่ถึง 0.1 ในทั้งสองทิศ
    \item \textbf{DoS} ได้ค่าต่ำถึงปานกลางในทั้งสองทิศ (0.12--0.28 และ 0.19--0.59) เทียบกับค่าเพดาน 0.79--0.98 และไม่เปลี่ยนเมื่อเปลี่ยนสัดส่วนการแบ่ง
    \item \textbf{DDoS} มีค่าเพดานเกือบสมบูรณ์ทั้งสองคอร์ปัส แต่ส่งผ่านในทิศ 2017\dirarrow{}2018 ได้สูงสุดเพียง 0.011 และในทิศ 2018\dirarrow{}2017 มีเพียง Logistic Regression ที่ได้เกิน 0.14 (0.518) DDoS จึงเป็นคลาสที่สูญเสียมากที่สุดเมื่อเทียบค่าเพดานกับค่าการส่งผ่าน
    \item \textbf{Bot} จำแนกได้ยากแม้ภายใน CICIDS2017 เอง (ค่าเพดาน 0.00--0.27) แต่ภายใน CSE-CIC-IDS2018 ได้เกือบ 1.0 ในทิศ 2018\dirarrow{}2017 มีเพียง LightGBM ที่ได้เกิน 0.05 (0.869)
    \item \textbf{Web Attack} ส่งผ่านไม่ได้ในทั้งสองทิศ ทั้งที่ค่าเพดานสูง
    \item \textbf{Benign} ได้ 0.88--0.95 ในทุกกรณี เนื่องจาก Macro F1-score เป็นค่าเฉลี่ยของ 7 คลาส คลาส Benign เพียงคลาสเดียวจึงคิดเป็นคะแนนประมาณ 0.13 ซึ่งเท่ากับค่าพื้นโดยประมาณ คะแนนส่วนที่สูงกว่าค่าพื้นจึงมาจากคลาสโจมตีเพียงไม่กี่คลาส ได้แก่ Brute Force, DoS และ Bot ในบางแบบจำลอง
    \item \textbf{Infiltration} รายงานไม่ได้อย่างน่าเชื่อถือ เพราะ CICIDS2017 มีตัวอย่างเพียง 36 ระเบียน (ชุดฝึก 25 ชุดทดสอบ 11) ค่า F1 ของทิศ 2018\dirarrow{}2017 จึงคำนวณจากระเบียนเพียง 11 ระเบียน ขณะที่แม้ภายใน CSE-CIC-IDS2018 เองแบบจำลองก็ได้ค่า F1 เพียง 0.00--0.22
\end{itemize}

\subsection{ความไวต่อสัดส่วนการแบ่งข้อมูล: กรณี SSH-Patator}
\label{sec:cd-ssh}

เมื่อเปลี่ยนสัดส่วนชุดฝึกจาก 60\% (protocol\_v2) เป็น 70\% (protocol\_v3) ผลภายใต้การแบ่งแบบสุ่มแทบไม่เปลี่ยน แต่ผลภายใต้การแบ่งตามลำดับการจับข้อมูลเปลี่ยนมาก สาเหตุอยู่ที่องค์ประกอบของคลาส Brute Force ใน CICIDS2017 ซึ่งประกอบด้วยเครื่องมือโจมตี 2 ตัวที่ทำงานต่อเนื่องกันในไฟล์เดียว คือ FTP-Patator แล้วตามด้วย SSH-Patator การแบ่งตามลำดับการจับข้อมูลตัดที่สัดส่วนของแถวในกลุ่ม ตำแหน่งที่ตัดจึงกำหนดว่าเครื่องมือใดอยู่ในชุดฝึก ดังตารางที่ \ref{tab:cd-ssh}

\begin{table}[ht]
    \caption{องค์ประกอบของคลาส Brute Force ใน CICIDS2017 ภายใต้การแบ่งตามลำดับการจับข้อมูล}
    \label{tab:cd-ssh}
    \small
    \begin{tabularx}{\linewidth}{@{}X r r r@{}}
        \hline
        \textbf{เครื่องมือโจมตี} & \textbf{ชุดฝึก (60\%)} & \textbf{ชุดฝึก (70\%)} & \textbf{ชุดทดสอบ (ทั้งสองแบบ)} \\
        \hline
        FTP-Patator & 5,490 & 5,931 & 0 \\
        SSH-Patator & \textbf{0} & \textbf{474} & 2,745 \\
        \hline
    \end{tabularx}
\end{table}

ที่ชุดฝึก 60\% แบบจำลองที่ฝึกจาก CICIDS2017 ไม่เคยเห็น SSH-Patator เลย ขณะที่ชุดทดสอบของคลาสนี้เป็น SSH-Patator ทั้งหมด 2,745 ระเบียน ที่ชุดฝึก 70\% รอยตัดเลื่อนเข้าไปในช่วงของ SSH-Patator ทำให้มี SSH-Patator 474 ระเบียนในชุดฝึก ซึ่งเป็นเครื่องมือเดียวกับ SSH-Bruteforce ของ CSE-CIC-IDS2018 ผลที่ตามมาแสดงในตารางที่ \ref{tab:cd-ssh-f1} และรูปที่ \ref{fig:cd-split-sensitivity}

\begin{table}[ht]
    \caption{ค่า F1 ของคลาส Brute Force และ Macro F1-score ในทิศ 2017\dirarrow{}2018 ภายใต้การแบ่งตามลำดับการจับข้อมูล}
    \label{tab:cd-ssh-f1}
    \small
    \begin{tabularx}{\linewidth}{@{}l *{4}{>{\raggedleft\arraybackslash}X}@{}}
        \hline
        \textbf{แบบจำลอง} & \textbf{Brute Force (60\%)} & \textbf{Brute Force (70\%)} & \textbf{Macro F1 (60\%)} & \textbf{Macro F1 (70\%)} \\
        \hline
        LightGBM & 0.0000 & 0.7729 & 0.2058 & 0.3132 \\
        XGBoost & 0.0000 & 0.7901 & 0.1712 & 0.3368 \\
        Stacking & 0.0000 & 0.9290 & 0.1715 & 0.3322 \\
        CatBoost / RF / MLP & 0.0000 & 0.0000 & 0.17--0.18 & 0.17--0.18 \\
        Logistic Regression & 0.0024 & 0.0000 & 0.1508 & 0.1496 \\
        \hline
    \end{tabularx}
\end{table}

\begin{figure}[ht]
    \centering
    \includegraphics[width=\linewidth]{fig_cd_split_sensitivity.png}
    \caption{ผลของการเปลี่ยนสัดส่วนชุดฝึกจาก 60\% เป็น 70\% ต่อค่า Macro F1-score (ซ้าย) และค่า F1 ของคลาส Brute Force (ขวา) ในทิศ 2017\dirarrow{}2018}
    \label{fig:cd-split-sensitivity}
\end{figure}

เนื่องจาก Macro F1-score เป็นค่าเฉลี่ยของ 7 คลาส การเปลี่ยนของ Macro F1-score จึงแยกได้เป็นผลรวมของการเปลี่ยนรายคลาส (คลาสละ ผลต่างของ F1 หารด้วย 7) ตารางที่ \ref{tab:cd-delta-class} แสดงว่าคลาสใดเป็นเจ้าของการเปลี่ยนแปลง

\begin{table}[ht]
    \caption{การเปลี่ยนของ Macro F1-score จากชุดฝึก 60\% เป็น 70\% แยกตามคลาส (การแบ่งตามลำดับการจับข้อมูล)}
    \label{tab:cd-delta-class}
    \footnotesize
    \begin{tabularx}{\linewidth}{@{}l *{6}{>{\raggedleft\arraybackslash}X}@{}}
        \hline
        \textbf{ค่าที่วัด} & \textbf{ผลต่าง Macro F1} & \textbf{Brute Force} & \textbf{Bot} & \textbf{DoS} & \textbf{DDoS} & \textbf{คลาสอื่น} \\
        \hline
        เพดาน 2017 (เฉลี่ย 7 แบบจำลอง) & +0.158 & +0.141 & 0.000 & 0.000 & 0.000 & +0.017 \\
        2017\dirarrow{}2018 (เฉลี่ย 7 แบบจำลอง) & +0.062 & +0.051 & +0.011 & 0.000 & 0.000 & 0.000 \\
        \quad LightGBM & +0.107 & +0.110 & 0.000 & -0.003 & 0.000 & 0.000 \\
        \quad XGBoost & +0.166 & +0.113 & +0.051 & +0.001 & 0.000 & +0.001 \\
        \quad Stacking & +0.161 & +0.133 & +0.027 & 0.000 & 0.000 & +0.001 \\
        2018\dirarrow{}2017 (เฉลี่ย 7 แบบจำลอง) & -0.006 & -0.001 & 0.000 & 0.000 & -0.004 & -0.001 \\
        \hline
    \end{tabularx}
\end{table}

\begin{itemize}
    \item ค่าเพดานของ CICIDS2017 ที่เพิ่มขึ้น 0.158 มาจาก Brute Force 0.141 (89\%) ส่วนที่เหลือเกือบทั้งหมดคือ Infiltration ซึ่งมีชุดทดสอบเพียง 11 ระเบียน
    \item ค่า Macro F1-score ของการส่งผ่านในทิศ 2017\dirarrow{}2018 ที่เพิ่มขึ้น 0.062 มาจาก Brute Force 82\% และ Bot 18\% การเปลี่ยนเกิดเฉพาะใน LightGBM, XGBoost และ Stacking ส่วนอีก 4 แบบจำลองเปลี่ยนไม่ถึง 0.002 ส่วน Bot ที่เพิ่มขึ้นใน XGBoost (F1 จาก 0.000 เป็น 0.358) และ Stacking (จาก 0.000 เป็น 0.187) ไม่ได้มาจากองค์ประกอบของชุดฝึก เมื่อแยกดูราย seed พบว่าค่า F1 ของ Bot ออกมาเพียงสองแบบ คือเกือบ 1.0 หรือ 0 ทั้งหมด XGBoost จำแนก Bot ได้ใน 0 จาก 5 รอบที่ชุดฝึก 60\% และ 2 จาก 5 รอบที่ชุดฝึก 70\% Stacking ได้ 0 และ 1 จาก 5 รอบ ส่วน LightGBM ได้ 1 จาก 5 รอบเท่ากันทั้งสองแบบ นอกจากนี้ Bot 195 ระเบียนที่เข้ามาในชุดฝึกเมื่อเลื่อนรอยตัดก็ไม่ได้ใกล้เคียง Bot ของ CSE-CIC-IDS2018 มากกว่า Bot ที่อยู่ในชุดฝึกอยู่แล้ว (มี Bot ของ CSE-CIC-IDS2018 เพียง 0.1\% ที่มีเพื่อนบ้านใกล้ที่สุดอยู่ในกลุ่มนี้ เทียบกับสัดส่วน 14.3\% ของกลุ่มนี้ในชุดฝึก) ค่าที่เพิ่มขึ้นจึงสะท้อนความไม่เสถียรของแบบจำลองระหว่างรอบการสุ่ม ไม่ใช่ผลของการเปลี่ยนสัดส่วนการแบ่ง
    \item ในทิศ 2018\dirarrow{}2017 ค่าเฉลี่ยแทบไม่เปลี่ยน แต่ใน Stacking คลาส DDoS ลดลงจาก 0.324 เป็น 0.132 ขณะที่ DoS เพิ่มจาก 0.384 เป็น 0.493 สองคลาสหักล้างกัน Macro F1-score จึงดูคงที่
    \item ความไม่สมมาตรภายใต้การแบ่งตามลำดับการจับข้อมูลลดลงจาก 33.6 จุด (ชุดฝึก 60\%) เป็น 13.8 จุด (ชุดฝึก 70\%) และค่าเพดานเฉลี่ยของ CICIDS2017 เพิ่มจาก 0.649 เป็น 0.807 ส่วนภายใต้การแบ่งแบบสุ่ม ความไม่สมมาตรเปลี่ยนจาก 4.3 เป็น 4.6 จุดเท่านั้น
\end{itemize}

ข้อสรุปของหัวข้อนี้คือ ภายใต้การแบ่งตามลำดับการจับข้อมูล องค์ประกอบของเครื่องมือโจมตีย่อยในชุดฝึกสามารถเปลี่ยนข้อสรุปเชิงปริมาณของการทดลองข้ามชุดข้อมูลได้มาก เพียงการเลื่อนรอยตัดให้ครอบคลุม SSH-Patator 474 ระเบียนทำให้ความไม่สมมาตรลดลงกว่าครึ่ง ขนาดของความไม่สมมาตรจึงควรรายงานเป็นช่วง (14--34 จุด) พร้อมระบุองค์ประกอบของชุดฝึก มิใช่ตัวเลขเดียว ขณะที่ทิศทางของความไม่สมมาตรยังคงเดิม

\subsection{ความแตกต่างเฉพาะชุดข้อมูล}
\label{sec:cd-av}

ผลการตรวจสอบเชิงปฏิปักษ์ตามวิธีในหัวข้อ \ref{sec:cd-av-method} แสดงในตารางที่ \ref{tab:cd-av} และรูปที่ \ref{fig:cd-av} การตรวจนี้ฝึกตัวจำแนกให้แยกคอร์ปัสบนแถวที่สุ่มจากแต่ละคลาสด้วยการแบ่งของตนเอง จึงไม่ขึ้นกับสัดส่วนการแบ่ง 60/10/30 หรือ 70/30 ของโพรโทคอลการส่งผ่าน

% แหล่งข้อมูล: round5_reliability_audit/data/diagnostics.json, within_dataset_control.json, tool_matched.json
\begin{table}[ht]
    \caption{ค่า AUC ของการแยกคอร์ปัส CICIDS2017 กับ CSE-CIC-IDS2018 รายคลาส (ค่าเฉลี่ย 5 รอบ)}
    \label{tab:cd-av}
    \footnotesize
    \begin{tabularx}{\linewidth}{@{}X r r r l r r@{}}
        \hline
        \textbf{คลาส} & \textbf{LightGBM} & \textbf{LR} & \textbf{Stump} & \textbf{ภายในคอร์ปัส} & \textbf{ซ้ำทุกค่า} & \textbf{ใกล้ซ้ำ} \\
        \hline
        Benign & 0.9990 & 0.9478 & 0.7089 & 0.6580--0.8112 & 0 & 0.10\% \\
        DDoS & 1.0000 & 1.0000 & 0.9917 & 0.9997--1.0000\textsuperscript{†} & 0 & 6.82\% \\
        DoS & 1.0000 & 0.9951 & 0.9370 & 0.9952--1.0000\textsuperscript{†} & 0 & 0.16\% \\
        Brute Force & 1.0000 & 0.9999 & 0.9950 & 0.8606--0.8993 & 0 & 0.20\% \\
        Bot & 1.0000 & 0.9990 & 0.9959 & 0.7584--0.9365 & 0 & 0.17\% \\
        Web Attack & 1.0000 & 1.0000 & 0.9992 & 0.7095--0.8404 & 0 & 2.75\% \\
        \hline
        \multicolumn{7}{@{}l}{\scriptsize ภายในคอร์ปัส: ช่วงของค่า AUC จากแบบจำลองทั้ง 7 รูปแบบ เมื่อแยกคนละไฟล์หรือคนละช่วงเวลาในคอร์ปัสเดียวกัน} \\
        \multicolumn{7}{@{}l}{\scriptsize \textsuperscript{†} ใน CSE-CIC-IDS2018 สองฝั่งของเส้นแบ่งเป็นเครื่องมือโจมตีคนละตัว ไม่ใช่เพียงคนละวัน} \\
    \end{tabularx}
\end{table}

\begin{figure}[ht]
    \centering
    \includegraphics[width=0.9\linewidth]{fig_cd_av.png}
    \caption{ช่วงของค่า AUC ในการแยกข้อมูลสองกลุ่ม ตั้งแต่กลุ่มควบคุมที่ไม่ควรแยกได้ไปจนถึงการแยกข้ามคอร์ปัส (กลุ่มควบคุมภายในคอร์ปัสไม่รวม DoS และ DDoS)}
    \label{fig:cd-av}
\end{figure}

\begin{itemize}
    \item การแยกข้ามคอร์ปัสได้ค่า AUC 0.9990--1.0000 ในทุกคลาสที่ทดสอบได้ และค่าแทบไม่ต่างกันระหว่าง 5 รอบ คลาส Infiltration ทดสอบไม่ได้ เพราะ CICIDS2017 มีเพียง 36 ระเบียน
    \item การสลับป้ายกำกับทำให้ค่า AUC ลดลงเหลือ 0.4928--0.5055 ในทุกแบบจำลองและทุกคลาส ส่วนกลุ่มควบคุมว่างได้ 0.49--0.53
    \item ภายในคอร์ปัสเดียวกัน คนละไฟล์หรือคนละช่วงเวลาแยกได้ 0.66--0.94 (ไม่รวม DoS และ DDoS) เมื่อตรึงเครื่องมือโจมตีให้เหมือนกันแล้วเปลี่ยนเฉพาะวัน ได้ 0.75--0.83 แต่เมื่อตรึงเครื่องมือโจมตีให้เหมือนกันแล้วเปลี่ยนเฉพาะคอร์ปัส ได้ 0.9999--1.0000
    \item Logistic Regression แยกได้ตั้งแต่ 0.9478 ขึ้นไป และ decision stump ซึ่งใช้กฎเดียวบนคอลัมน์เดียวแยกคลาสโจมตีได้ 0.937--0.999 (Benign ได้ 0.709)
    \item ไม่มีระเบียนที่ซ้ำกันทุกค่าคร่อมชุดฝึกและชุดทดสอบเลยในการสุ่มทั้ง 30 ครั้ง ระเบียนใกล้ซ้ำอยู่ใกล้ระดับบังเอิญ (0.10--0.20\%) ยกเว้น DDoS (6.82\%) และ Web Attack (2.75\%) แต่ decision stump ซึ่งจำแถวไม่ได้ก็ยังแยกสองคลาสนี้ได้ 0.99 ขึ้นไป
\end{itemize}

สรุปได้ว่าสองคอร์ปัสแยกออกจากกันได้เกือบสมบูรณ์ และความต่างนี้แรงกว่าความต่างระหว่างวันหรือไฟล์ภายในคอร์ปัสเดียวกันอย่างชัดเจน หลักฐานที่บ่งชี้ว่าค่า AUC ที่สูงไม่ได้มาจากการรั่วไหลแบบตรง ๆ มีสามข้อ คือ การสลับป้ายกำกับให้ค่าระดับสุ่ม ไม่มีระเบียนซ้ำคร่อมชุดฝึกและชุดทดสอบ และแบบจำลองที่มีความจุต่ำมากอย่าง decision stump ก็แยกได้ ความต่างดังกล่าวจึงอยู่ในค่าของคุณลักษณะเอง และกระจายอยู่ในคุณลักษณะจำนวนมาก โดยเฉพาะกลุ่มขนาดแพ็กเก็ต ซึ่งเป็นกลุ่มเดียวกับที่ใช้จำแนกการโจมตี (สหสัมพันธ์อันดับระหว่างความสามารถแยกคอร์ปัสกับ mutual information ต่อคลาสโจมตีเท่ากับ +0.72 ใน CICIDS2017 และ +0.77 ใน CSE-CIC-IDS2018)

\subsection{การวิเคราะห์สาเหตุของการทำนายพลาดด้วย SHAP}
\label{sec:cd-shap}

การทดสอบสมมติฐานตามหัวข้อ \ref{sec:cd-shap-method} บนการแบ่ง 70/30 ให้ผลดังนี้

\begin{itemize}
    \item การเทียบภายในคอร์ปัสปลายทาง (กลุ่ม F หลัก 13 ตัว) มีกรณีที่จำนวนข้อมูลพอให้ตัดสินได้ 11 กรณีจาก 4 แบบจำลองและสองทิศ (DoS 8 กรณี และอีก 3 กรณีคือ LightGBM Brute Force ทิศ 2017\dirarrow{}2018, LightGBM Bot และ XGBoost DDoS ทิศ 2018\dirarrow{}2017) \textbf{ไม่ผ่านเกณฑ์ทั้ง 11 กรณี} ค่า \textit{G}\textsubscript{within} ติดลบ 9 จาก 11 กรณี และค่า \textit{p} ต่ำสุดคือ 0.174
    \item กรณี DoS ทิศ 2017\dirarrow{}2018 ตอบคำถามได้ชัดที่สุด เมื่อเทียบกับระเบียนที่ทำนายถูกในคอร์ปัสต้นทาง ส่วนต่างเป็นบวกใน LightGBM, XGBoost และ CatBoost แต่เมื่อเทียบกับระเบียนที่ทำนายถูกในคอร์ปัสปลายทางเดียวกัน \textit{G}\textsubscript{within} ติดลบในทั้ง 4 แบบจำลอง (-0.07 ถึง -0.48) และติดลบทั้ง 20 ค่าราย seed แสดงว่าส่วนต่างบวกเดิมมาจากการที่ข้อมูลมาจากต่างคอร์ปัส ไม่ได้มาจากการทำนายผิด
    \item การเทียบกับกลุ่มคุณลักษณะสุ่มในทิศ 2017\dirarrow{}2018 (เทียบกับระเบียนที่ทำนายถูกในคอร์ปัสต้นทาง) มีกรณีที่นับได้ 17 กรณี ไม่ผ่านทั้งหมด ค่า \textit{p} ต่ำสุด 0.063 (CatBoost DoS)
\end{itemize}

ข้อค้นพบรองที่ต้องรายงานอย่างระมัดระวังคือ เมื่อใช้กลุ่มคุณลักษณะที่เลือกด้วยเกณฑ์ SHAP อย่างเดียว (26 ตัว) มี 2 กรณีที่ผ่านเกณฑ์รายคลาส ได้แก่ XGBoost DoS (\textit{G} = +0.228, \textit{p} = 0.0045) และ CatBoost DDoS (\textit{G} = +0.246, \textit{p} = 0.0090) ทั้งสองไม่ใช่กลุ่มหลักที่ใช้ตัดสิน และไม่ทำให้ผลรวมของแบบจำลองนั้นเปลี่ยน เมื่อตรวจซ้ำแบบภายในคอร์ปัสปลายทาง XGBoost DoS ได้ \textit{G}\textsubscript{within} = +0.073 (\textit{p} = 0.21) ซึ่งไม่ผ่าน ส่วน CatBoost DDoS ตรวจไม่ได้ เพราะแบบจำลองทำนาย DDoS ใน CSE-CIC-IDS2018 ถูก 0 ระเบียน นอกจากนี้ การวิเคราะห์รอบเสริมซึ่งระบุเกณฑ์ไว้ก่อนการวิเคราะห์เช่นกัน ได้ใช้กลุ่มคุณลักษณะ 13 ตัวที่ค่าของคลาสโจมตีนั้นเลื่อนมากที่สุดระหว่างสองคอร์ปัสแทนกลุ่ม F ผลรวมตามเกณฑ์เป็นแบบบางส่วน โดยคลาสที่ชัดที่สุดคือ DoS ซึ่งผ่านเกณฑ์ใน 3 จาก 4 แบบจำลอง อย่างไรก็ตาม รอบนี้เทียบกับระเบียนที่ทำนายถูกในคอร์ปัสต้นทาง ซึ่งการเทียบภายในคอร์ปัสปลายทางข้างต้นแสดงแล้วว่ายังปนผลของการที่ข้อมูลมาจากต่างคอร์ปัส ผลนี้จึงบ่งชี้ได้เพียงว่าการเลื่อนค่าระดับคลาสอาจมีส่วนต่อการทำนายพลาด และต้องตรวจซ้ำด้วยการเทียบภายในคอร์ปัสปลายทางก่อนนำไปสรุป

ข้อสรุปที่ใช้ได้อย่างปลอดภัยคือ \emph{ไม่พบหลักฐานว่ากลุ่มคุณลักษณะที่บ่งชี้ชุดข้อมูลซึ่งกำหนดไว้ก่อนการวิเคราะห์เป็นสาเหตุโดยตรงของการทำนายระเบียนโจมตีเป็น Benign หลังการส่งผ่านข้ามชุดข้อมูล}

%======================================================================
\clearpage
\section{สรุปข้อค้นพบสำคัญ}
\label{sec:key-findings}

ข้อค้นพบด้านล่างใช้ protocol\_v3 (ชุดฝึก 70\% / ชุดทดสอบ 30\%) เป็นหลัก ตัวเลขของชุดฝึก 60\% ใช้เพื่อเปรียบเทียบความไวเท่านั้น

\begin{enumerate}
    \item \textbf{การทดลองหลัก} CatBoost ได้ Macro F1-score สูงสุดบน CICIDS2017 แบบ 9 คลาสที่ 0.8873 ตรวจจับการโจมตีได้ 98.40\% ที่ FPR 0.00065 และเป็นอันดับหนึ่งในทั้งสามเกณฑ์การจัดอันดับ
    \item \textbf{ค่าการส่งผ่านสูงสุด} ภายใต้การแบ่งตามลำดับการจับข้อมูล ทิศ 2017\dirarrow{}2018 ได้สูงสุด 0.3368 (XGBoost) และทิศ 2018\dirarrow{}2017 ได้สูงสุด 0.4435 (LightGBM) เทียบกับค่าเพดานของแบบจำลองเดียวกัน 0.8589 และ 0.8767 และค่าพื้นประมาณ 0.13
    \item \textbf{ค่าดิบใกล้ค่าพื้น} ทิศ 2017\dirarrow{}2018 มี 4 จาก 7 แบบจำลองที่ได้เพียง 0.15--0.18 (ค่าพื้น 0.1337)
    \item \textbf{ความไม่สมมาตร} ทิศ 2018\dirarrow{}2017 กู้คืนได้มากกว่า โดยค่า Recovery เฉลี่ย 28.2\% เทียบกับ 14.4\% (ต่างกัน 13.8 จุด, 6 จาก 7 แบบจำลอง) ภายใต้การแบ่งตามลำดับการจับข้อมูล และ 15.2\% เทียบกับ 10.6\% (4.6 จุด, 4 จาก 7 แบบจำลอง) ภายใต้การแบ่งแบบสุ่ม
    \item \textbf{การแบ่งแบบสุ่มให้ค่าเพดานของ CICIDS2017 สูงกว่าการแบ่งตามลำดับการจับข้อมูล} ค่าเพดานเฉลี่ย 0.9065 (สุ่ม) เทียบกับ 0.8073 (ตามลำดับการจับข้อมูล) ขณะที่ CSE-CIC-IDS2018 ต่างกันเพียง 0.012
    \item \textbf{ความไวต่อสัดส่วนการแบ่ง} SSH-Patator ในชุดฝึกของ CICIDS2017 เพิ่มจาก 0 เป็น 474 ระเบียนเมื่อเปลี่ยนจาก 60\% เป็น 70\% ทำให้ความไม่สมมาตรลดจาก 33.6 เป็น 13.8 จุด และค่าเพดานเฉลี่ยของ CICIDS2017 เพิ่มจาก 0.649 เป็น 0.807
    \item \textbf{Brute Force} ค่า F1 ในทิศ 2017\dirarrow{}2018 เพิ่มจาก 0 เป็น 0.77 (LightGBM), 0.79 (XGBoost) และ 0.93 (Stacking) และอธิบายการเพิ่มของ Macro F1-score เฉลี่ยได้ 82\% (ที่เหลือ 18\% มาจาก Bot)
    \item \textbf{คลาสที่ส่งผ่านไม่ได้} DDoS ในทิศ 2017\dirarrow{}2018 ได้สูงสุดเพียง 0.011 และ Web Attack ได้สูงสุด 0.002 และ 0.018 ในสองทิศ ทั้งที่ค่าเพดานเฉลี่ยอยู่ที่ 0.99 และ 0.78--0.80
    \item \textbf{ความแตกต่างเฉพาะชุดข้อมูล} AUC ของการแยกคอร์ปัส 0.9990--1.0000 ทุกคลาส การสลับป้ายกำกับได้ 0.49--0.51 ภายในคอร์ปัสเดียวกันได้ 0.66--0.94 Logistic Regression ได้ตั้งแต่ 0.948 และไม่มีระเบียนซ้ำคร่อมชุดฝึกและชุดทดสอบ
    \item \textbf{SHAP} การเทียบภายในคอร์ปัสปลายทางผ่าน 0 จาก 11 กรณี (\textit{p} ต่ำสุด 0.174) และ \textit{G}\textsubscript{within} ของ DoS ทิศ 2017\dirarrow{}2018 ติดลบในทั้ง 4 แบบจำลอง
\end{enumerate}

    \include{chapter5}

    \clearpage

    \def\inbib{1}
    \addcontentsline{toc}{chapter}{\bibname}
    \bibliographystyle{IEEEtran-kmitl}
    \makeatletter
    \EveryShipout{
        \ifdim\pagetotal>\pagegoal
            \if\inbib1
                \aftergroup\@cont@heading
            \else
            \fi
        \fi
    }
    \makeatother

    \bibliography{reference}
    \def\inbib{0}

    \startappendix
    \include{appendix}
    \chapter{เมทริกซ์ความสับสนของการทดลองหลักบน CICIDS2017}
\label{appendix:main-confusion-matrices}

ภาคผนวกนี้แสดงเมทริกซ์ความสับสนแบบจำนวนนับของอีก 6 แบบจำลองบนชุดทดสอบเดียวกับ CatBoost ในรูปที่ \ref{fig:main-cm-catboost} แถวแสดงคลาสจริงและคอลัมน์แสดงคลาสที่ทำนาย

\begin{figure}[!htbp]
    \centering
    \includegraphics[width=0.92\linewidth]{lightgbm_confusion_matrix.png}
    \caption{เมทริกซ์ความสับสนของ LightGBM จากการทดลองหลักบนชุดทดสอบ CICIDS2017}
    \label{fig:main-cm-lightgbm}
\end{figure}

\clearpage

\begin{figure}[!htbp]
    \centering
    \includegraphics[width=0.92\linewidth]{xgboost_confusion_matrix.png}
    \caption{เมทริกซ์ความสับสนของ XGBoost จากการทดลองหลักบนชุดทดสอบ CICIDS2017}
    \label{fig:main-cm-xgboost}
\end{figure}

\clearpage

\begin{figure}[!htbp]
    \centering
    \includegraphics[width=0.92\linewidth]{random_forest_confusion_matrix.png}
    \caption{เมทริกซ์ความสับสนของ Random Forest จากการทดลองหลักบนชุดทดสอบ CICIDS2017}
    \label{fig:main-cm-random-forest}
\end{figure}

\clearpage

\begin{figure}[!htbp]
    \centering
    \includegraphics[width=0.92\linewidth]{stacking_confusion_matrix.png}
    \caption{เมทริกซ์ความสับสนของ Stacking Ensemble จากการทดลองหลักบนชุดทดสอบ CICIDS2017}
    \label{fig:main-cm-stacking}
\end{figure}

\clearpage

\begin{figure}[!htbp]
    \centering
    \includegraphics[width=0.92\linewidth]{logistic_regression_confusion_matrix.png}
    \caption{เมทริกซ์ความสับสนของ Logistic Regression จากการทดลองหลักบนชุดทดสอบ CICIDS2017}
    \label{fig:main-cm-logistic-regression}
\end{figure}

\clearpage

\begin{figure}[!htbp]
    \centering
    \includegraphics[width=0.92\linewidth]{mlp_confusion_matrix.png}
    \caption{เมทริกซ์ความสับสนของ Multi-Layer Perceptron จากการทดลองหลักบนชุดทดสอบ CICIDS2017}
    \label{fig:main-cm-mlp}
\end{figure}


    \makeauthorbio

\end{document}
